%=====================================================================
% INSTALACOES ELETRICAS PREDIAIS - LIGACAO, MONTAGEM E VERIFICACAO
% Parte 2 de 2 do material predial
% Atualizacao: agosto de 2026
%=====================================================================
\documentclass[11pt,aspectratio=169]{beamer}

\usepackage[utf8]{inputenc}
\usepackage[T1]{fontenc}
\usepackage[brazil]{babel}
\usepackage{lmodern}
\usepackage{amsmath,amssymb,bm}
\usepackage{siunitx}
\usepackage{booktabs,array,tabularx,multirow}
\usepackage[most]{tcolorbox}
\usepackage{tikz}
\usepackage{pgfplots}
\usepackage[european,american currents]{circuitikz}
\usepackage{ragged2e}
\usepackage{microtype}
\usepackage{xcolor}
\usetikzlibrary{arrows.meta,positioning,calc,fit,shapes.geometric,patterns,decorations.pathmorphing,decorations.markings,matrix}
\pgfplotsset{compat=1.18}
\sisetup{locale=DE,per-mode=symbol,detect-all}

\newcommand{\sen}{\operatorname{sen}}
\newcommand{\tg}{\operatorname{tg}}
\newcommand{\fase}[1]{\textcolor{corPrim}{\textbf{#1}}}
\newcommand{\alerta}[1]{\textcolor{corAt}{\textbf{#1}}}
\newcommand{\Disciplina}{}
\newcommand{\TituloCurto}{Montagem}
%=====================================================================
% ESTILO DA CASA -- TEMPLATE BEAMER (padrao da colecao)
% Arquivo central de identidade visual para todos os materiais.
% Azul-petroleo + ambar | tcolorbox | TikZ/pgfplots/circuitikz
%=====================================================================

% Dependencias minimas do proprio estilo.
\RequirePackage{xcolor}
\RequirePackage{tikz}
\usetikzlibrary{calc}
\RequirePackage{tcolorbox}
\tcbuselibrary{skins,breakable}
\RequirePackage{listings}
\RequirePackage{xparse}
\RequirePackage{etoolbox}

\makeatletter

%====================== PALETA OFICIAL ================================
\definecolor{corPrim}{RGB}{13,48,82}      % azul-petroleo
\definecolor{corAccent}{RGB}{222,138,28}  % ambar
\definecolor{corDef}{RGB}{31,111,168}     % azul medio
\definecolor{corNorma}{RGB}{0,121,107}    % verde-petroleo
\definecolor{corEx}{RGB}{46,125,50}       % verde
\definecolor{corFix}{RGB}{204,102,0}      % laranja
\definecolor{corAt}{RGB}{176,58,46}       % vermelho
\definecolor{codBg}{RGB}{252,250,245}
\definecolor{codKw}{RGB}{13,48,82}
\definecolor{codStr}{RGB}{46,125,50}
\definecolor{codCom}{RGB}{120,130,138}

% Compatibilidade com nomes usados historicamente nos materiais.
\colorlet{Petroleo}{corPrim}
\colorlet{PetroleoEscuro}{corPrim!82!black}
\colorlet{PetroleoMedio}{corDef!82!corPrim}
\colorlet{PetroleoClaro}{corDef!7}
\colorlet{Ambar}{corAccent}
\colorlet{AmbarEscuro}{corAccent!72!black}
\colorlet{AmbarClaro}{corAccent!16}
\colorlet{AzulPetroleo}{corPrim}
\colorlet{AzulPetroleoEscuro}{corPrim!82!black}
\colorlet{AzulPetroleoClaro}{corDef!7}
\colorlet{AzulMedio}{corDef}
\colorlet{AzulClaro}{corDef!7}
\colorlet{CinzaEscuro}{black!82}
\colorlet{CinzaClaro}{black!5}
\colorlet{CinzaSuave}{black!3}
\colorlet{CinzaMedio}{black!55}
\colorlet{CinzaTexto}{black!82}
\colorlet{CinzaBox}{black!3}
\colorlet{Verde}{corEx}
\colorlet{Vermelho}{corAt}
\colorlet{Turquesa}{corNorma!70!corDef}
\colorlet{Roxo}{corDef!58!corAt}
\colorlet{ifmagreen}{corPrim}
\colorlet{ifmagreen2}{corDef}
\colorlet{ifmalight}{corDef!7}
\colorlet{accent}{corAccent}
\colorlet{warn}{corAt}
\colorlet{ok}{corEx}
\colorlet{bluex}{corDef}
\colorlet{grayx}{black!58}

% Variantes em minusculas presentes em materiais antigos.
\colorlet{azulPetroleo}{corPrim}
\colorlet{azulPetroleoEscuro}{corPrim!82!black}
\colorlet{azulPetroleoClaro}{corDef!7}
\colorlet{azulMedio}{corDef}
\colorlet{azulClaro}{corDef!7}
\colorlet{azulEscuro}{corPrim!82!black}
\colorlet{azul}{corPrim}
\colorlet{azul2}{corDef}
\colorlet{ambar}{corAccent}
\colorlet{ambarClaro}{corAccent!16}
\colorlet{cinza}{black!5}
\colorlet{cinzaClaro}{black!5}
\colorlet{cinzaMedio}{black!55}
\colorlet{cinzaEscuro}{black!82}
\colorlet{cinzaTexto}{black!82}
\colorlet{verde}{corEx}
\colorlet{verdeclaro}{corEx!10}
\colorlet{vermelho}{corAt}
\colorlet{vermelhoclaro}{corAt!9}
\colorlet{laranja}{corFix}
\colorlet{roxo}{corDef!58!corAt}

%====================== TEMA BEAMER ==================================
\usetheme{default}
\usefonttheme{professionalfonts}
\setbeamercolor{normal text}{fg=black!85,bg=white}
\setbeamercolor{structure}{fg=corPrim}
\setbeamercolor{frametitle}{bg=corPrim,fg=white}
\setbeamercolor{footsec}{bg=corPrim,fg=white!92!corPrim}
\setbeamercolor{alerted text}{fg=corAccent!85!black}
\setbeamerfont{frametitle}{series=\bfseries,size=\large}
\setbeamerfont{footline}{size=\tiny}
\setbeamertemplate{navigation symbols}{}
\setbeamertemplate{caption}{\raggedright\insertcaption\par}
\setbeamersize{text margin left=0.55cm,text margin right=0.55cm}

\setbeamertemplate{itemize item}{\raise0.35ex\hbox{\color{corAccent}\rule{0.62ex}{0.62ex}}}
\setbeamertemplate{itemize subitem}{\color{corDef}\raise0.5pt\hbox{--}}
\setbeamertemplate{itemize subsubitem}{\color{corPrim}\tiny$\circ$}
\setbeamertemplate{enumerate item}{\color{corPrim}\textbf{\insertenumlabel.}}
\setbeamertemplate{section in toc}{%
  \leavevmode\textcolor{corAccent}{\textbf{\inserttocsectionnumber.}}~\inserttocsection\par\vskip1.5pt}
\setbeamertemplate{subsection in toc}{%
  \leavevmode\leftskip=0.95cm\textcolor{corDef}{\textbullet}~\inserttocsubsection\par}

%====================== TITULO DO SLIDE ==============================
\setbeamertemplate{frametitle}{%
  \nointerlineskip%
  \begin{beamercolorbox}[wd=\paperwidth,ht=2.7ex,dp=1.3ex,leftskip=0.55cm,rightskip=0.55cm]{frametitle}%
    \usebeamerfont{frametitle}\insertframetitle%
  \end{beamercolorbox}%
  {\color{corAccent}\rule{\paperwidth}{1.3pt}}\par\vskip2pt%
}

%====================== RODAPE DUPLO ================================
\newcommand{\CasaRodapeEsq}{%
  \ifcsname Disciplina\endcsname\Disciplina\else\insertshorttitle\fi}
\newcommand{\CasaRodapeDir}{%
  \ifcsname TituloCurto\endcsname\TituloCurto\else\insertshorttitle\fi}
\newcommand{\CasaMetaCapa}{%
  \ifcsname Disciplina\endcsname
    \if\relax\detokenize\expandafter{\Disciplina}\relax
    \else
      \Disciplina\ \textbullet\ \insertdate
    \fi
  \else
    \insertdate
  \fi}
\setbeamertemplate{footline}{%
  \leavevmode%
  \hbox{%
    \begin{beamercolorbox}[wd=.5\paperwidth,ht=2.4ex,dp=1ex,leftskip=0.5cm]{footsec}%
      \usebeamerfont{footline}\CasaRodapeEsq%
    \end{beamercolorbox}%
    \begin{beamercolorbox}[wd=.5\paperwidth,ht=2.4ex,dp=1ex,rightskip=0.5cm]{footsec}%
      \hfill\usebeamerfont{footline}\CasaRodapeDir\quad\textbar\quad\insertframenumber%
    \end{beamercolorbox}%
  }\vskip0pt%
}

%====================== DIVISORIAS DE SECAO =========================
\newcommand{\CasaSecNum}{\ifnum\value{section}<10 0\fi\arabic{section}}
\newcommand{\CasaDivisoriaManual}[2]{%
  \begin{frame}[plain,noframenumbering]
    \makebox[0pt][l]{%
      \raisebox{0pt}[0pt][0pt]{%
        \begin{tikzpicture}[remember picture,overlay]
          \fill[corPrim] (current page.south west) rectangle (current page.north east);
          \if\relax\detokenize{#1}\relax\else
            \node[anchor=west,text=corAccent,font=\bfseries\fontsize{40}{40}\selectfont]
              at ([xshift=1.6cm,yshift=1.1cm]current page.west) {#1};
          \fi
          \fill[corAccent] ([xshift=1.7cm,yshift=0.35cm]current page.west)
            rectangle ++(2.4,0.09);
          \node[anchor=west,text=white,font=\bfseries\LARGE,text width=0.75\paperwidth]
            at ([xshift=1.6cm,yshift=-0.5cm]current page.west) {#2};
        \end{tikzpicture}%
      }%
    }%
  \end{frame}%
}
\AtBeginSection[]{\CasaDivisoriaManual{\CasaSecNum}{\insertsection}}

% Alguns materiais antigos usam \secslide{titulo}; outros, \secslide{numero}{titulo}.
\@ifundefined{secslide}
  {\NewDocumentCommand{\secslide}{m g}{\IfNoValueTF{#2}{\CasaDivisoriaManual{}{#1}}{\CasaDivisoriaManual{#1}{#2}}}}
  {\RenewDocumentCommand{\secslide}{m g}{\IfNoValueTF{#2}{\CasaDivisoriaManual{}{#1}}{\CasaDivisoriaManual{#1}{#2}}}}
\@ifundefined{Divisoria}
  {\NewDocumentCommand{\Divisoria}{m g}{\IfNoValueTF{#2}{\CasaDivisoriaManual{}{#1}}{\CasaDivisoriaManual{#1}{#2}}}}
  {\RenewDocumentCommand{\Divisoria}{m g}{\IfNoValueTF{#2}{\CasaDivisoriaManual{}{#1}}{\CasaDivisoriaManual{#1}{#2}}}}

%====================== CAPA ========================================
\defbeamertemplate*{title page}{estilo da casa}{%
  \vspace{0.8cm}%
  {\color{white}\bfseries\huge\inserttitle\par}%
  \vspace{2mm}%
  {\color{corAccent}\rule{3.2cm}{2.2pt}\par}%
  \vspace{4mm}%
  {\color{white!85}\large\insertsubtitle\par}%
  \vspace{8mm}%
  {\color{white}\normalsize\insertauthor\par}%
  \vspace{1mm}%
  {\color{white!70}\small\CasaMetaCapa\par}%
}
\setbeamertemplate{title page}[estilo da casa]

%====================== CAIXAS =======================================
% Materiais antigos usam convencoes diferentes para o argumento opcional
% das caixas. Para preservar a semantica, caixas existentes nao sao
% redefinidas; apenas sua paleta historica e recolorida acima.
\tcbset{
  caixabase/.style={
    enhanced,sharp corners=downhill,boxrule=0.6pt,left=2mm,right=2mm,
    top=1.2mm,bottom=1.2mm,fonttitle=\bfseries\small,
    attach boxed title to top left={xshift=3mm,yshift=-2.4mm},
    boxed title style={sharp corners,boxrule=0pt},
    before skip=3pt,after skip=3pt
  }
}
\@ifundefined{caixaDef}{\newtcolorbox{caixaDef}[1][]{caixabase,colback=corDef!5,colframe=corDef,coltitle=white,colbacktitle=corDef,title=Definição,#1}}{}
\@ifundefined{caixaNorma}{\newtcolorbox{caixaNorma}[1][]{caixabase,colback=corNorma!5,colframe=corNorma,coltitle=white,colbacktitle=corNorma,title=Norma / Método,#1}}{}
\@ifundefined{caixaEx}{\newtcolorbox{caixaEx}[1][]{caixabase,colback=corEx!5,colframe=corEx,coltitle=white,colbacktitle=corEx,title=Exemplo resolvido,#1}}{}
\@ifundefined{caixaFix}{\newtcolorbox{caixaFix}[1][]{caixabase,colback=corFix!5,colframe=corFix,coltitle=white,colbacktitle=corFix,title=Fixação,#1}}{}
\@ifundefined{caixaAt}{\newtcolorbox{caixaAt}[1][]{caixabase,colback=corAt!5,colframe=corAt,coltitle=white,colbacktitle=corAt,title=Atenção,#1}}{}
\@ifundefined{caixaForm}{\newtcolorbox{caixaForm}[1][]{caixabase,colback=corPrim!4,colframe=corPrim,coltitle=white,colbacktitle=corPrim,title=Formulário,#1}}{}
\@ifundefined{caixaObs}{\newtcolorbox{caixaObs}[1][]{caixabase,colback=black!3,colframe=corPrim,coltitle=white,colbacktitle=corPrim,title=Observação,#1}}{}
\@ifundefined{caixaAlerta}{\newtcolorbox{caixaAlerta}[1][]{caixabase,colback=corAt!5,colframe=corAt,coltitle=white,colbacktitle=corAt,title=Atenção,#1}}{}
\@ifundefined{caixaProjeto}{\newtcolorbox{caixaProjeto}[1][]{caixabase,colback=corNorma!5,colframe=corNorma,coltitle=white,colbacktitle=corNorma,title=Projeto prático,#1}}{}
\@ifundefined{caixaProj}{\newtcolorbox{caixaProj}[1][]{caixabase,colback=corNorma!5,colframe=corNorma,coltitle=white,colbacktitle=corNorma,title=Projeto prático,#1}}{}

%====================== CODIGO =======================================
\lstdefinestyle{codCasa}{
  basicstyle=\ttfamily\scriptsize,
  keywordstyle=\color{codKw}\bfseries,
  stringstyle=\color{codStr},commentstyle=\color{codCom}\itshape,
  numbers=left,numberstyle=\tiny\color{codCom},numbersep=5pt,
  frame=single,rulecolor=\color{corPrim!25},backgroundcolor=\color{codBg},
  breaklines=true,showstringspaces=false,tabsize=2,columns=fullflexible
}
\lstdefinestyle{codC}{style=codCasa,language=C}
\lstdefinestyle{codPy}{style=codCasa,language=Python}
\lstdefinestyle{py}{style=codCasa,language=Python}

%====================== EXTENSOES MODULARES ==========================
% Se existir <jobname>_antes_laboratorios.tex, o arquivo e carregado uma
% unica vez imediatamente antes da secao de laboratorios. Isso permite
% ampliar um material sem duplicar o arquivo principal.
\newif\ifCasaExtensaoAntesLabsInserida
\CasaExtensaoAntesLabsInseridafalse
\AtBeginDocument{%
  \IfFileExists{\jobname_antes_laboratorios.tex}{%
    \let\CasaSectionOriginal\section
    \renewcommand{\section}[1]{%
      \ifCasaExtensaoAntesLabsInserida\else
        \ifstrequal{##1}{Laboratórios, projetos e comissionamento}{%
          \global\CasaExtensaoAntesLabsInseridatrue
          \input{\jobname_antes_laboratorios.tex}%
        }{}%
      \fi
      \CasaSectionOriginal{##1}%
    }%
  }{}%
}

\makeatother

%=====================================================================
% SIMBOLOS ELETRICOS TIKZ — DEFINICOES SEM SLIDES
% Carregado no preambulo dos materiais que precisam desenhar plantas.
%=====================================================================

\providecommand{\SimLuz}[2]{\begin{scope}[shift={(#1,#2)}]
\draw[corAccent,thick] (0,0) circle (0.14);
\draw[corAccent,thick] (-0.10,-0.10)--(0.10,0.10);
\draw[corAccent,thick] (-0.10,0.10)--(0.10,-0.10);
\end{scope}}
\providecommand{\SimTUG}[2]{\begin{scope}[shift={(#1,#2)}]
\draw[corPrim,thick] (0,0) circle (0.13);
\draw[corPrim,thick] (-0.16,0)--(0.16,0);
\end{scope}}
\providecommand{\SimTUGM}[2]{\begin{scope}[shift={(#1,#2)}]
\draw[corPrim,thick] (0,0) circle (0.13);
\draw[corPrim,thick] (-0.16,0)--(0.16,0);
\draw[corPrim,thick] (-0.16,0.065)--(0.16,0.065);
\end{scope}}
\providecommand{\SimTUE}[3]{\begin{scope}[shift={(#1,#2)}]
\draw[corFix,thick] (0,0) circle (0.16);
\node[font=\tiny\bfseries,corFix] at (0,0) {#3};
\end{scope}}
\providecommand{\SimInterruptor}[3]{\begin{scope}[shift={(#1,#2)}]
\draw[corNorma,thick] (0,0) circle (0.11);
\node[font=\tiny\bfseries,corNorma,anchor=west] at (0.14,0) {#3};
\end{scope}}
\providecommand{\SimQD}[3]{\begin{scope}[shift={(#1,#2)}]
\draw[corPrim,fill=corDef!10,thick,rounded corners=1pt] (-0.32,-0.20) rectangle (0.32,0.20);
\node[font=\tiny\bfseries] at (0,0) {#3};
\end{scope}}
\providecommand{\SimDados}[2]{\begin{scope}[shift={(#1,#2)}]
\draw[corDef,thick] (0,0.15)--(-0.14,-0.11)--(0.14,-0.11)--cycle;
\node[font=\tiny\bfseries,corDef] at (0,-0.01) {D};
\end{scope}}
\providecommand{\SimMotor}[3]{\begin{scope}[shift={(#1,#2)}]
\draw[corAt,thick] (0,0) circle (0.20);
\node[font=\tiny\bfseries,corAt] at (0,0) {#3};
\end{scope}}
\providecommand{\SimVFD}[3]{\begin{scope}[shift={(#1,#2)}]
\draw[corDef,fill=corDef!8,thick,rounded corners=1pt] (-0.28,-0.18) rectangle (0.28,0.18);
\node[font=\tiny\bfseries,corDef] at (0,0) {#3};
\end{scope}}
\providecommand{\SimTomadaTri}[2]{\begin{scope}[shift={(#1,#2)}]
\draw[corFix,thick] (0,0) circle (0.16);
\node[font=\tiny\bfseries,corFix] at (0,0) {$3\sim$};
\end{scope}}
\providecommand{\SimEmerg}[2]{\begin{scope}[shift={(#1,#2)}]
\draw[corAt,fill=corAt!12,thick] (0,0) circle (0.16);
\node[font=\tiny\bfseries,corAt] at (0,0) {!};
\end{scope}}

\providecommand{\RotaL}[4]{\draw[corAccent!80!black,dashed,thick] (#1,#2)--(#3,#4);}
\providecommand{\RotaT}[4]{\draw[corPrim!75,dashed,thick] (#1,#2)--(#3,#4);}
\providecommand{\RotaF}[4]{\draw[corFix!80,dashed,thick] (#1,#2)--(#3,#4);}

\providecommand{\CotaH}[4]{%
\draw[thin] (#1,#3)--(#1,#3+0.28);
\draw[thin] (#2,#3)--(#2,#3+0.28);
\draw[{Latex}-{Latex},thin] (#1,#3+0.20)--(#2,#3+0.20)
node[midway,fill=white,inner sep=1pt,font=\tiny]{#4};}
\providecommand{\CotaV}[4]{%
\draw[thin] (#3,#1)--(#3+0.28,#1);
\draw[thin] (#3,#2)--(#3+0.28,#2);
\draw[{Latex}-{Latex},thin] (#3+0.20,#1)--(#3+0.20,#2)
node[midway,fill=white,inner sep=1pt,font=\tiny,rotate=90]{#4};}


\title{Instalações Elétricas Prediais}
\subtitle{Parte 2 — Ligação, montagem, execução e verificação}
\author{Francisco}
\institute{Do esquema elétrico à instalação verificada}
\date{}

\begin{document}

{\setbeamercolor{background canvas}{bg=corPrim}\setbeamertemplate{footline}{}
\begin{frame}[plain,noframenumbering]\titlepage\end{frame}}

\begin{frame}{O que este Beamer entrega}
\small
\begin{columns}[T,onlytextwidth]
\column{0.49\textwidth}
\begin{caixaProjeto}[title=Ligação e montagem]
\begin{itemize}
\item leitura de planta, unifilar e esquemas funcionais;
\item ferramentas, materiais, caixas e eletrodutos;
\item lançamento, identificação e terminação de condutores;
\item iluminação, tomadas e cargas específicas;
\item montagem completa do quadro de distribuição.
\end{itemize}
\end{caixaProjeto}
\column{0.49\textwidth}
\begin{caixaNorma}[title=Segurança e entrega]
\begin{itemize}
\item desenergização, bloqueio e organização do trabalho;
\item continuidade, isolação, polaridade e testes de DR;
\item inspeção, energização controlada e comissionamento;
\item falhas típicas e diagnóstico sistemático;
\item prática integrada com lista de materiais e registros.
\end{itemize}
\end{caixaNorma}
\end{columns}
\end{frame}

\section{Segurança, NR-10 e manutenção}

\begin{frame}{Desenergização: lógica de segurança}
\small
\centering
\begin{tikzpicture}[node distance=0.30cm,every node/.style={font=\scriptsize}]
\tikzset{b/.style={draw=corPrim,fill=corDef!7,rounded corners,minimum width=2.4cm,minimum height=0.65cm,align=center}}
\node[b] (a) {Seccionar};
\node[b,right=of a] (b) {Impedir\\reenergização};
\node[b,right=of b] (c) {Constatar ausência\\de tensão};
\node[b,below=of c] (d) {Aterramento\\temporário, se aplicável};
\node[b,left=of d] (e) {Proteger contra\\partes energizadas};
\node[b,left=of e] (f) {Sinalizar o impedimento\\de reenergização};
\foreach \u/\v in {a/b,b/c,c/d,d/e,e/f}{\draw[-{Stealth},thick,corPrim] (\u)--(\v);}
\end{tikzpicture}
\vspace{3mm}
\begin{caixaObs}
A sequência real deve seguir procedimento aprovado, análise de risco e requisitos aplicáveis. O slide mostra a lógica, não substitui procedimento operacional.
\end{caixaObs}
\end{frame}

\begin{frame}{LOTO: bloqueio e etiquetagem}
\small
\begin{columns}[T,onlytextwidth]
\column{0.50\textwidth}
\begin{itemize}
\item Identificar todas as fontes de energia.
\item Seccionar e bloquear cada fonte aplicável.
\item Dissipar/neutralizar energias armazenadas.
\item Etiquetar responsável e condição.
\item Verificar estado seguro antes do serviço.
\item Restabelecer somente após procedimento de liberação.
\end{itemize}
\column{0.48\textwidth}
\centering
\begin{tikzpicture}[scale=0.86,every node/.style={font=\scriptsize}]
\draw[thick,fill=corAt!8,rounded corners=3pt] (0,0) rectangle (4.8,3.2);
\node[font=\bfseries\large,corAt] at (2.4,2.65) {BLOQUEADO};
\draw[thick] (1.4,0.7) rectangle (3.4,1.8);
\draw[thick] (2.4,1.8) arc (180:0:0.7 and 0.85);
\node at (2.4,0.35) {não operar};
\end{tikzpicture}
\end{columns}
\end{frame}

\begin{frame}{Prontuário e documentação de segurança}
\small
\begin{columns}[T,onlytextwidth]
\column{0.49\textwidth}
\begin{itemize}
\item Diagramas atualizados.
\item Procedimentos e instruções técnicas.
\item Inspeções e medições do sistema de proteção.
\item Registros de treinamento/autorização conforme aplicável.
\item Especificação de EPC/EPI.
\end{itemize}
\column{0.49\textwidth}
\begin{itemize}
\item Estudos de risco pertinentes.
\item Relatórios de ensaio e manutenção.
\item Histórico de modificações.
\item Dados de curto-circuito e ajustes quando necessários à segurança.
\item Plano de emergência aplicável.
\end{itemize}
\end{columns}
\end{frame}

\begin{frame}{Termografia: exemplo de interpretação}
\small
\begin{columns}[T,onlytextwidth]
\column{0.48\textwidth}
\begin{itemize}
\item Comparar componentes equivalentes sob carga semelhante.
\item Investigar corrente, aperto, oxidação, dimensionamento e ventilação.
\item Temperatura absoluta isolada não explica a causa.
\item Registrar carga, ambiente e condições da inspeção.
\end{itemize}
\column{0.50\textwidth}
\begin{tikzpicture}
\begin{axis}[ybar,width=5.8cm,height=3.8cm,bar width=16pt,ylabel={$T$ (°C)},symbolic x coords={A,B,C,N},xtick=data,ymin=20,ymax=100,nodes near coords,grid=major]
\addplot coordinates {(A,48)(B,51)(C,82)(N,55)};
\end{axis}
\end{tikzpicture}
\end{columns}
\begin{caixaAt}
Fase C muito mais quente é um \textbf{indício}; a ação correta é diagnosticar antes de apenas ``reapertar tudo''.
\end{caixaAt}
\end{frame}

%=====================================================================

%=====================================================================
% MONTAGEM PREDIAL - SEGURANCA, PLANEJAMENTO E FERRAMENTAS
%=====================================================================

\section{Planejamento seguro, ferramentas e materiais}

\begin{frame}{Montagem começa antes de abrir a primeira caixa}
\centering
\begin{tikzpicture}[
  b/.style={draw=corPrim,fill=corDef!7,rounded corners,minimum width=2.15cm,minimum height=0.78cm,align=center,font=\scriptsize},
  a/.style={-{Stealth},thick,corPrim}]
\node[b](proj)at(-5.0,0){Projeto\\liberado};
\node[b](risco)at(-2.5,0){Análise de risco\\e procedimento};
\node[b](mat)at(0,0){Materiais e\\instrumentos};
\node[b](exec)at(2.5,0){Execução\\controlada};
\node[b](ensaio)at(5.0,0){Inspeção, ensaio\\e registro};
\draw[a](proj)--(risco);\draw[a](risco)--(mat);\draw[a](mat)--(exec);\draw[a](exec)--(ensaio);
\end{tikzpicture}
\vspace{5mm}
\begin{caixaAt}
Este material é didático. Intervenções em instalações elétricas exigem trabalhadores qualificados, capacitados e autorizados conforme as responsabilidades e condições aplicáveis.
\end{caixaAt}
\end{frame}

\begin{frame}{Responsabilidades precisam estar definidas}
\small
\begin{tabularx}{\textwidth}{>{\bfseries}p{2.9cm}X X}
\toprule
Papel & Antes do serviço & Durante e depois \\
\midrule
Responsável técnico & projeto, especificações e critérios de aceitação & decisões técnicas, alterações e documentação \\
Supervisor & equipe, recursos, análise de risco e sequência & controle da execução e liberação das etapas \\
Executor autorizado & conferir tarefa, ferramentas e condição segura & executar, comunicar desvios e registrar resultados \\
Responsável por ensaios & método, instrumento e limites & medir, interpretar e emitir registro \\
Cliente/operação & acesso, cargas e restrições de uso & aceite, operação e manutenção futura \\
\bottomrule
\end{tabularx}
\end{frame}

\begin{frame}{Ordem de serviço: transformar projeto em tarefa controlada}
\small
\begin{enumerate}
\item identificar local, circuito, origem e limites físicos da atividade;
\item anexar desenho, revisão e lista de materiais aplicáveis;
\item registrar perigos elétricos e não elétricos;
\item definir desenergização, impedimento de reenergização e liberação;
\item atribuir responsáveis, equipe e comunicação;
\item listar instrumentos, EPIs, EPCs e calibração necessária;
\item estabelecer inspeções intermediárias e critérios de aceite;
\item fechar alterações, ensaios, fotos e as built.
\end{enumerate}
\begin{caixaObs}
Uma ordem de serviço genérica não controla risco nem qualidade. A tarefa precisa ter origem, destino, revisão e critério de término.
\end{caixaObs}
\end{frame}

\begin{frame}{Análise de risco: eletricidade não é o único perigo}
\small
\begin{tabularx}{\textwidth}{>{\bfseries}p{2.9cm}X X}
\toprule
Perigo & Exposição típica & Controle planejado \\
\midrule
Choque/arco & quadro, condutor existente, retorno e geração & desenergização, teste, bloqueio, barreiras \\
Altura & luminárias, fachadas e entrada aérea & acesso adequado, NR-35 quando aplicável \\
Perfuração & parede com tubulação ou circuito oculto & compatibilização, detecção e delimitação \\
Poeira/projeção & cortes, rasgos e furações & proteção coletiva, ocular e respiratória \\
Esforço/queda & bobinas, escadas, ferramentas e entulho & organização, movimentação e área isolada \\
Ambiente & umidade, calor, chuva e espaço reduzido & suspender, adaptar método ou controlar condição \\
\bottomrule
\end{tabularx}
\end{frame}

\begin{frame}{Ferramentas: escolher pelo processo, não pelo hábito}
\small
\begin{tabularx}{\textwidth}{>{\bfseries}p{2.6cm}X X}
\toprule
Etapa & Ferramentas principais & Controle \\
\midrule
Marcação & trena, nível, esquadro, linha e detector & cotas e interferências \\
Infraestrutura & furadeira, serra-copo, curvador, guia e escareador & diâmetro, raio e acabamento \\
Condutores & guia, lubrificante compatível, desenrolador e puxador & tração e torção \\
Preparação & decapador, alicate de corte e desencapador de cabo & não ferir a alma \\
Terminação & prensa-terminal e torquímetro/chave de torque & matriz e torque corretos \\
Identificação & impressora/etiqueta e marcador permanente & padrão único e legível \\
\bottomrule
\end{tabularx}
\end{frame}

\begin{frame}{Instrumentos: cada um responde a uma pergunta}
\small
\begin{tabularx}{\textwidth}{>{\bfseries}p{2.9cm}X X}
\toprule
Instrumento & Pergunta & Limite de uso \\
\midrule
Detector bipolar & há tensão e qual a faixa? & verificar antes/depois em fonte conhecida \\
Multímetro TRMS & tensão, resistência, frequência? & categoria, ponta e fusível adequados \\
Alicate amperímetro & corrente sem abrir o circuito? & abraçar o condutor correto, não o cabo completo \\
Megômetro & isolação suporta tensão de ensaio? & desconectar eletrônica sensível \\
Terrômetro & resistência do sistema/eletrodo? & método e geometria compatíveis \\
Testador de DR & corrente/tempo de atuação? & interpretar conforme dispositivo e circuito \\
Analisador de rede & harmônicas, potência e eventos? & configuração, janela e referência de tensão \\
\bottomrule
\end{tabularx}
\end{frame}

\begin{frame}{Ferramenta isolada não autoriza trabalho energizado}
\small
\begin{columns}[T,onlytextwidth]
\column{0.49\textwidth}
\begin{caixaAt}[title=Erro perigoso]
Confundir cabo isolado, cabo de ferramenta revestido ou luva comum com proteção certificada para o risco elétrico.
\end{caixaAt}
\column{0.49\textwidth}
\begin{caixaNorma}[title=Uso correto]
\begin{itemize}
\item ferramenta compatível com a tarefa e tensão;
\item inspeção visual antes do uso;
\item armazenamento que preserve a isolação;
\item retirada imediata quando houver dano;
\item procedimento de trabalho prevalece sobre o acessório.
\end{itemize}
\end{caixaNorma}
\end{columns}
\begin{caixaObs}
Trabalho energizado é exceção tecnicamente justificada, com análise e controles específicos. A regra de montagem predial é preparar e executar desenergizado.
\end{caixaObs}
\end{frame}

\begin{frame}{EPI e EPC: barreiras complementares}
\small
\begin{tabularx}{\textwidth}{>{\bfseries}p{2.7cm}X X}
\toprule
Grupo & Exemplos & Verificar \\
\midrule
Proteção coletiva & bloqueio, barreira, sinalização, cobertura e aterramento temporário quando aplicável & instalação, estabilidade e responsabilidade \\
Cabeça/face & capacete e proteção facial compatível & classe, estado e ajuste \\
Mãos & luvas adequadas ao risco e cobertura mecânica quando necessária & ensaio, validade e perfuração \\
Corpo & vestimenta compatível com o risco térmico e ambiente & fechamento, conservação e contaminação \\
Olhos/respiração & óculos, viseira e proteção respiratória & poeira, projeção e ventilação \\
Altura & sistema de acesso e proteção contra queda & inspeção e ancoragem \\
\bottomrule
\end{tabularx}
\begin{caixaAt}
EPI não corrige circuito energizado sem necessidade, bloqueio ausente ou procedimento incompleto.
\end{caixaAt}
\end{frame}

\begin{frame}{Recebimento: rejeitar material antes de instalar}
\small
\begin{columns}[T,onlytextwidth]
\column{0.49\textwidth}
\begin{itemize}
\item conferir fabricante, modelo e certificação aplicável;
\item comparar seção, tensão, temperatura e classe do cabo;
\item verificar curva, polos, corrente e capacidade do disjuntor;
\item confirmar tipo e sensibilidade do DR;
\item conferir $U_c$, $U_p$, $I_n/I_{max}$ do DPS;
\item inspecionar trincas, oxidação, umidade e lacres.
\end{itemize}
\column{0.49\textwidth}
\begin{caixaProjeto}[title=Rastreabilidade]
Registrar:
\begin{itemize}
\item lote e nota de recebimento;
\item ficha técnica usada;
\item quantidade aceita/rejeitada;
\item local de armazenamento;
\item substituições aprovadas.
\end{itemize}
\end{caixaProjeto}
\end{columns}
\end{frame}

\begin{frame}{Organização do posto reduz falha de montagem}
\small
\begin{columns}[T,onlytextwidth]
\column{0.52\textwidth}
\begin{tikzpicture}[scale=0.82]
\draw[fill=black!3] (0,0) rectangle (6.2,3.8);
\draw[fill=corDef!8,draw=corPrim] (0.35,2.45) rectangle (2.25,3.35);
\node[font=\scriptsize,align=center] at (1.3,2.9){componentes\\conferidos};
\draw[fill=corNorma!8,draw=corNorma] (2.55,2.45) rectangle (4.45,3.35);
\node[font=\scriptsize,align=center] at (3.5,2.9){ferramentas\\inspecionadas};
\draw[fill=corAccent!10,draw=corAccent] (4.75,2.45) rectangle (5.85,3.35);
\node[font=\scriptsize,align=center] at (5.3,2.9){resíduos};
\draw[very thick] (0.45,0.75)--(5.75,0.75);
\node[font=\scriptsize] at (3.1,1.15){bancada limpa e iluminada};
\node[font=\tiny] at (3.1,0.35){documento e circuito identificados};
\end{tikzpicture}
\column{0.46\textwidth}
\begin{itemize}
\item separar materiais por circuito;
\item manter pontas e ferramentas protegidas;
\item evitar sobras dentro de caixas e quadros;
\item controlar parafusos, terminais e etiquetas;
\item proteger equipamentos contra poeira e impacto;
\item registrar etapa concluída antes de fechar.
\end{itemize}
\end{columns}
\end{frame}

\begin{frame}{Liberação da etapa de montagem}
\small
\begin{enumerate}
\item projeto e revisão corretos estão disponíveis no local;
\item circuito e fontes possíveis foram identificados;
\item desenergização, bloqueio e teste foram concluídos;
\item equipe, responsabilidades e comunicação estão definidos;
\item ferramentas, instrumentos e EPIs/EPCs foram inspecionados;
\item material confere com a especificação;
\item ambiente permite execução segura e limpa;
\item critérios de inspeção e registro estão definidos.
\end{enumerate}
\begin{caixaAt}
Qualquer resposta negativa interrompe a liberação. Corrigir a condição antes de iniciar a execução.
\end{caixaAt}
\end{frame}

%=====================================================================
% MONTAGEM PREDIAL - INFRAESTRUTURA, CONDUTORES E CONEXOES
%=====================================================================

\section{Infraestrutura, condutores e conexões}

\begin{frame}{Infraestrutura define se o projeto pode ser montado}
\centering
\begin{tikzpicture}[
  b/.style={draw=corPrim,fill=corDef!7,rounded corners,minimum width=1.85cm,minimum height=0.72cm,align=center,font=\scriptsize},
  a/.style={-{Stealth},thick,corPrim}]
\node[b](m)at(-5.2,0){Marcação\\e cotas};
\node[b](c)at(-3.1,0){Caixas e\\quadros};
\node[b](e)at(-1.0,0){Eletrodutos\\e rotas};
\node[b](l)at(1.1,0){Limpeza e\\guia};
\node[b](p)at(3.2,0){Puxamento\\dos cabos};
\node[b](t)at(5.3,0){Terminação\\e inspeção};
\draw[a](m)--(c);\draw[a](c)--(e);\draw[a](e)--(l);\draw[a](l)--(p);\draw[a](p)--(t);
\end{tikzpicture}
\vspace{5mm}
\begin{caixaAt}
Não fechar parede, forro ou piso antes de conferir continuidade da rota, diâmetro, segregação, acessibilidade e correspondência com o projeto.
\end{caixaAt}
\end{frame}

\begin{frame}{Marcação: transferir o desenho para a obra}
\small
\begin{columns}[T,onlytextwidth]
\column{0.48\textwidth}
\begin{tikzpicture}[scale=0.78]
\draw[very thick] (0,0) rectangle (6.1,4.4);
\draw[fill=black!8] (0,0) rectangle (6.1,0.25);
\draw[dashed,corAccent,thick] (0,1.25)--(6.1,1.25);
\draw[dashed,corDef,thick] (0,2.55)--(6.1,2.55);
\node[draw=corPrim,minimum size=0.5cm] at (1.2,1.25){T};
\node[draw=corNorma,minimum size=0.5cm] at (3.0,1.25){S};
\node[draw=corPrim,minimum width=1.1cm,minimum height=0.65cm] at (4.9,2.55){QD};
\draw[{Latex}-{Latex},thin] (5.75,0.25)--(5.75,1.25)node[midway,left,font=\tiny]{cota do projeto};
\end{tikzpicture}
\column{0.50\textwidth}
\begin{itemize}
\item adotar uma referência de nível única;
\item medir pelo eixo/face definidos no desenho;
\item conferir abertura de portas e posição do mobiliário;
\item compatibilizar com hidráulica, gás, estrutura e dados;
\item marcar circuito e tipo de caixa, não apenas um ponto;
\item registrar alteração antes de rasgar ou perfurar.
\end{itemize}
\end{columns}
\begin{caixaObs}
Alturas são definidas no projeto e nas exigências de acessibilidade/uso. Não existe uma única cota prática válida para todo ambiente.
\end{caixaObs}
\end{frame}

\begin{frame}{Caixas: volume, função e acesso}
\small
\begin{tabularx}{\textwidth}{>{\bfseries}p{2.6cm}X X}
\toprule
Tipo & Uso & Conferir \\
\midrule
4x2 / 4x4 & tomadas, interruptores e derivações & quantidade de módulos, profundidade e condutores \\
Octogonal & ponto de iluminação e suporte previsto & carga mecânica e fixação independente quando exigida \\
Passagem & mudança de direção e puxamento & dimensões, tampa e acesso permanente \\
Quadro embutido & proteção e distribuição & módulos, reserva, barramentos e dissipação \\
Caixa de piso & energia/dados em áreas abertas & vedação, resistência e segregação \\
Externa & áreas úmidas ou expostas & grau IP, prensa-cabos, dreno e material \\
\bottomrule
\end{tabularx}
\begin{caixaAt}
Caixa não é depósito de sobra de cabo. Acomodação deve preservar raio de curvatura, isolação, terminais e acesso aos parafusos.
\end{caixaAt}
\end{frame}

\begin{frame}{Eletrodutos: continuidade mecânica e raio de curvatura}
\small
\begin{columns}[T,onlytextwidth]
\column{0.51\textwidth}
\begin{tikzpicture}[scale=0.82]
\draw[corPrim,line width=8pt] (0,0)--(2.4,0) arc (-90:0:1.35) --(3.75,3.1);
\draw[white,line width=3pt] (0,0)--(2.4,0) arc (-90:0:1.35) --(3.75,3.1);
\draw[dashed,corAccent,thick] (2.4,0)--(3.75,1.35);
\node[font=\scriptsize,corAccent] at (2.45,1.25){raio};
\node[draw=corPrim,fill=corDef!7,minimum size=0.65cm] at (0,0){C};
\node[draw=corPrim,fill=corDef!7,minimum size=0.65cm] at (3.75,3.1){C};
\end{tikzpicture}
\column{0.47\textwidth}
\begin{itemize}
\item manter diâmetro e seção útil ao longo da rota;
\item eliminar rebarbas e bordas cortantes;
\item respeitar raio do eletroduto e dos cabos;
\item fixar para evitar deslocamento na concretagem;
\item usar buchas, arruelas e luvas compatíveis;
\item inserir caixas quando o puxamento exigir.
\end{itemize}
\end{columns}
\end{frame}

\begin{frame}{Antes do puxamento: provar e limpar a rota}
\small
\begin{enumerate}
\item conferir origem, destino e identificação do eletroduto;
\item passar guia com ponteira que não agrida a tubulação;
\item verificar obstrução, esmagamento e água acumulada;
\item limpar sem introduzir material incompatível;
\item medir o comprimento real e acrescentar folgas previstas;
\item confirmar ocupação e esforço de puxamento;
\item proteger bordas de entrada e saída;
\item posicionar bobinas para desenrolar sem torção.
\end{enumerate}
\begin{caixaAt}
Não usar o novo cabo como ferramenta para desobstruir eletroduto. Corrigir a infraestrutura antes de lançar os condutores.
\end{caixaAt}
\end{frame}

\begin{frame}{Ocupação: conferir com o diâmetro externo real}
\small
Para três ou mais cabos no eletroduto:
\[
\eta=\frac{\sum A_{cabos}}{A_{interna}}\leq0{,}40
\]
\begin{caixaEx}[title={Rota com 3 fases, neutro e PE}]
Cinco cabos com diâmetro externo de \SI{4.8}{mm}:
\[
\sum A=5\frac{\pi(4{,}8)^2}{4}=\SI{90.5}{mm^2}
\]
Para eletroduto com diâmetro interno de \SI{18.5}{mm}:
\[
\eta=\frac{90{,}5}{\pi(18{,}5)^2/4}=33{,}7\%
\]
A ocupação atende, mas ainda é preciso verificar curvas, distância e esforço de lançamento.
\end{caixaEx}
\end{frame}

\begin{frame}{Desenrolar corretamente evita torção permanente}
\centering
\begin{tikzpicture}[scale=0.9]
\draw[very thick] (-4,0) circle (1.15);
\draw[very thick] (-4,0) circle (0.35);
\draw[very thick] (-5.35,-1.35)--(-2.65,-1.35);
\draw[-{Stealth},thick,corNorma] (-2.8,0.7)--(-1.4,0.7);
\draw[corPrim,thick] (-2.85,0.25).. controls (-0.8,0.5) and (0.8,0.0)..(3.8,0.25);
\node[font=\scriptsize,corNorma] at (-2.1,1.05){giro da bobina};
\node[font=\scriptsize] at (-4,-1.75){desenrolador};
\node[draw=corPrim,fill=corDef!7,minimum width=1.7cm,minimum height=0.65cm] at (4.6,0.25){eletroduto};
\draw[corAt,thick,dashed] (-4,-2.45) circle (0.8);
\draw[corAt,thick] (-4.55,-3.0)--(-3.45,-1.9);
\node[font=\scriptsize,corAt,anchor=west] at (-3.0,-2.45){não retirar em espiras pelo lado};
\end{tikzpicture}
\begin{caixaObs}
Bobina apoiada em eixo ou desenrolador gira enquanto o cabo sai tangencialmente. Isso reduz laços, torção, dano e esforço de puxamento.
\end{caixaObs}
\end{frame}

\begin{frame}{Puxamento: distribuir esforço e proteger a isolação}
\small
\begin{columns}[T,onlytextwidth]
\column{0.50\textwidth}
\begin{enumerate}
\item agrupar e escalonar as pontas;
\item aplicar malha/puxador apropriado;
\item proteger a transição com fita compatível;
\item usar lubrificante aprovado para cabo/eletroduto;
\item manter comunicação entre alimentação e tração;
\item parar diante de esforço anormal;
\item deixar folga apenas onde prevista.
\end{enumerate}
\column{0.48\textwidth}
\begin{tikzpicture}[scale=0.82]
\foreach \y/\c in {0.75/corPrim,0.25/corAccent,-0.25/corNorma,-0.75/black!65}{
  \draw[\c,line width=3pt] (-3,\y)--(-0.4,0.18*\y);
}
\draw[fill=black!10] (-0.65,-0.55)--(0.65,-0.25)--(0.65,0.25)--(-0.65,0.55)--cycle;
\draw[very thick] (0.65,-0.25)--(2.4,-0.25)--(2.4,0.25)--(0.65,0.25);
\draw[-{Stealth},very thick,corPrim] (2.4,0)--(3.45,0);
\node[font=\tiny] at (0,-0.9){cabeça escalonada};
\end{tikzpicture}
\end{columns}
\end{frame}

\begin{frame}{Identificação: reconhecer o circuito nas duas pontas}
\small
\begin{tabularx}{\textwidth}{>{\bfseries}p{3.1cm}X X}
\toprule
Condutor & Identificação funcional & Regra de montagem \\
\midrule
PE & verde-amarelo & uso exclusivo como proteção \\
Neutro & azul-claro & manter coerência em toda a instalação \\
Fases & cores distintas do PE e neutro & repetir padrão no quadro e nas caixas \\
Retorno & cor diferente da fase permanente & etiquetar quando houver múltiplos retornos \\
Comando e\newline viajantes & cores e marcadores definidos & identificar origem, destino e função \\
Circuito & etiqueta em ambas as extremidades & código igual ao quadro e à planta \\
\bottomrule
\end{tabularx}
\begin{caixaAt}
Fita colorida pontual não transforma um condutor reservado a PE ou neutro em fase. A identificação deve respeitar as restrições normativas.
\end{caixaAt}
\end{frame}

\begin{frame}{Decapagem: remover isolação sem reduzir a seção}
\centering
\begin{tikzpicture}[scale=0.92]
\draw[corPrim,line width=10pt] (-5,1)--(-1.2,1);
\draw[corAccent,line width=4pt] (-1.2,1)--(0.8,1);
\node[font=\scriptsize] at (-3.1,1.65){isolação preservada};
\node[font=\scriptsize] at (-0.2,1.65){comprimento conforme borne};
\draw[corPrim,line width=10pt] (1.5,1)--(4.4,1);
\draw[corAccent,line width=4pt] (4.4,1)--(5.5,1);
\draw[corAt,very thick] (4.55,0.75)--(4.8,1.25);
\draw[corAt,very thick] (4.85,0.75)--(5.1,1.25);
\node[font=\scriptsize,corAt] at (3.7,0.15){alma cortada: rejeitar e refazer};
\end{tikzpicture}
\vspace{3mm}
\begin{itemize}
\item ajustar ferramenta à seção e à classe do condutor;
\item não marcar, torcer, estanhar ou cortar fios da alma;
\item respeitar o comprimento indicado pelo borne/terminal;
\item refazer a ponta sempre que houver dano ou cobre oxidado.
\end{itemize}
\end{frame}

\begin{frame}{Terminal: compatibilizar condutor, borne e ferramenta}
\small
\begin{tabularx}{\textwidth}{>{\bfseries}p{2.5cm}X X}
\toprule
Solução & Aplicação típica & Verificar \\
\midrule
Ponta direta & borne projetado para a classe do condutor & fios contidos e comprimento correto \\
Ilhós tubular & borne que aceita terminal comprimido & seção, comprimento e matriz \\
Olhal/garfo & parafuso ou barramento apropriado & diâmetro, material e pressão de contato \\
Luva de emenda & prolongamento/derivação autorizada & faixa de seções e isolamento posterior \\
Conector de alavanca/mola & derivação e manutenção rápida & certificação, seção, classe e corrente \\
Bimetálico & transição Cu-Al quando prevista & preparação, composto e torque \\
\bottomrule
\end{tabularx}
\begin{caixaAt}
Não colocar dois condutores em borne projetado para um, nem usar terminal maior ``compensado'' por esmagamento.
\end{caixaAt}
\end{frame}

\begin{frame}{Crimpagem: o resultado depende da matriz}
\centering
\begin{tikzpicture}[scale=0.86]
\node[draw=corPrim,fill=corDef!7,rounded corners,minimum width=2.0cm,minimum height=0.72cm](a)at(-4.5,0){seção e terminal};
\node[draw=corPrim,fill=corDef!7,rounded corners,minimum width=2.0cm,minimum height=0.72cm](b)at(-1.5,0){matriz correta};
\node[draw=corPrim,fill=corDef!7,rounded corners,minimum width=2.0cm,minimum height=0.72cm](c)at(1.5,0){compressão total};
\node[draw=corNorma,fill=corNorma!7,rounded corners,minimum width=2.0cm,minimum height=0.72cm](d)at(4.5,0){inspeção e tração};
\draw[-{Stealth},thick,corPrim](a)--(b);\draw[-{Stealth},thick,corPrim](b)--(c);\draw[-{Stealth},thick,corPrim](c)--(d);
\end{tikzpicture}
\vspace{4mm}
\begin{columns}[T,onlytextwidth]
\column{0.49\textwidth}
\begin{caixaAt}[title=Subcrimpagem]
Baixa pressão deixa vazios, aumenta resistência e permite soltura.
\end{caixaAt}
\column{0.49\textwidth}
\begin{caixaAt}[title=Sobrecrimpagem]
Matriz inadequada pode trincar o terminal ou reduzir a seção do condutor.
\end{caixaAt}
\end{columns}
\end{frame}

\begin{frame}{Torque: aperto correto é um valor, não uma sensação}
\small
\begin{columns}[T,onlytextwidth]
\column{0.48\textwidth}
\begin{tikzpicture}[scale=0.86]
\draw[fill=corDef!8,draw=corPrim,very thick] (-1.5,-0.55) rectangle (1.5,0.55);
\draw[fill=black!15] (-0.35,0.55) rectangle (0.35,1.25);
\draw[very thick] (-2.9,1.25)--(2.9,1.25);
\draw[corAccent,line width=4pt] (-1.2,0)--(1.2,0);
\node[font=\scriptsize] at (0,-1.05){borne com condutor preparado};
\draw[-{Stealth},thick,corNorma] (2.3,2.0) arc (30:-55:1.0);
\node[font=\scriptsize,corNorma] at (2.65,2.15){torque};
\end{tikzpicture}
\column{0.50\textwidth}
\begin{itemize}
\item obter o torque na ficha ou no próprio equipamento;
\item usar ferramenta adequada e dentro da faixa;
\item aplicar no parafuso correto, sem extensão improvisada;
\item marcar/registrar quando o procedimento exigir;
\item não reapertar periodicamente sem recomendação técnica;
\item investigar aquecimento, não apenas ``dar mais aperto''.
\end{itemize}
\end{columns}
\end{frame}

\begin{frame}{Derivações: manter conexão acessível e inspecionável}
\small
\begin{columns}[T,onlytextwidth]
\column{0.49\textwidth}
\begin{caixaNorma}[title=Conexão aceitável]
\begin{itemize}
\item conector adequado às seções e materiais;
\item dentro de caixa com tampa acessível;
\item isolação restabelecida e sem esforço mecânico;
\item identificação preservada;
\item volume suficiente para acomodação.
\end{itemize}
\end{caixaNorma}
\column{0.49\textwidth}
\begin{caixaAt}[title=Conexão a rejeitar]
\begin{itemize}
\item torção manual coberta por fita como solução final;
\item emenda enterrada no reboco ou dentro do eletroduto;
\item cobre exposto, fios soltos ou dois cabos no mesmo borne;
\item material incompatível ou sem faixa declarada;
\item cabo tracionando diretamente a conexão.
\end{itemize}
\end{caixaAt}
\end{columns}
\end{frame}

\begin{frame}{Inspeção antes de fechar caixas e rotas}
\small
\begin{enumerate}
\item caixas estão niveladas, firmes, limpas e na cota correta;
\item eletrodutos têm continuidade, acabamento e segregação;
\item cabos correspondem ao circuito, seção e cores do projeto;
\item isolação não sofreu corte, esmagamento ou tração anormal;
\item terminações usam terminal, matriz e torque corretos;
\item derivações estão acessíveis, isoladas e acomodadas;
\item folgas permitem manutenção sem excesso solto;
\item fotos e desvios foram registrados para o as built.
\end{enumerate}
\begin{caixaProjeto}
Fechar acabamento é um ponto de não retorno caro. A inspeção intermediária evita quebrar parede para corrigir falha escondida.
\end{caixaProjeto}
\end{frame}


% A nova sequencia substitui integralmente os dois capitulos antigos.
% Garante que nos TikZ com texto multilinha tenham alinhamento valido,
% inclusive quando o proprio desenho redefine apenas a fonte dos nos.
\tikzset{execute at begin picture={\tikzset{every node/.append style={align=center}}}}
%=====================================================================
% ESQUEMAS REVISADOS - ILUMINACAO E COMANDOS
%=====================================================================
\section{Esquemas de ligação — iluminação e comandos}

\definecolor{LigF}{RGB}{35,35,35}
\definecolor{LigN}{RGB}{0,95,175}
\definecolor{LigPE}{RGB}{35,145,65}
\definecolor{LigR}{RGB}{190,45,45}
\definecolor{LigVum}{RGB}{222,138,28}
\definecolor{LigVdois}{RGB}{112,70,155}
\tikzset{
  schbox/.style={draw=corPrim,fill=corDef!5,rounded corners=1.5pt,minimum width=2cm,minimum height=.8cm,align=center,font=\scriptsize},
  schsw/.style={draw=corNorma,fill=corNorma!5,rounded corners=1.5pt,minimum width=2cm,minimum height=1.2cm,align=center,font=\scriptsize},
  schload/.style={draw=corAccent!85!black,fill=corAccent!7,circle,minimum size=.9cm,align=center,font=\scriptsize},
  schterm/.style={circle,draw=corPrim,fill=white,inner sep=0pt,minimum size=4.8pt},
  schwire/.style={line width=1.05pt}
}

\begin{frame}{Padrão gráfico — o esquema precisa mostrar função e borne}
\small
\begin{columns}[T,onlytextwidth]
\column{.56\textwidth}
\begin{tikzpicture}[scale=.8,every node/.style={font=\scriptsize}]
\node[schbox](q)at(-3.7,1.1){QD / proteção};
\node[schbox](cx)at(0,1.1){caixa / passagem};
\node[schsw](s)at(0,-1.2){dispositivo\\bornes};
\node[schload](l)at(3.7,1.1){carga};
\draw[LigF,schwire](q.east)--(cx.west)node[midway,above]{F};
\draw[LigF,schwire](cx.south)--(s.north)node[midway,right]{F};
\draw[LigR,schwire](s.east)--(2.0,-1.2)--(2.0,1.1)--(l.west)node[pos=.45,right]{R};
\draw[LigN,schwire](-4.7,2.35)--(3.7,2.35)--(l.north)node[pos=.15,above]{N};
\draw[LigPE,schwire](-4.7,-2.25)--(4.15,-2.25)--(4.15,1.1)--(l.east)node[pos=.15,below]{PE};
\end{tikzpicture}
\column{.42\textwidth}
\begin{caixaProjeto}[title=Um bom desenho responde]
\begin{itemize}
\item origem e proteção;
\item borne de entrada e saída;
\item condutor permanente/comandado;
\item destino de N e PE;
\item trecho físico de cada fio.
\end{itemize}
\end{caixaProjeto}
\end{columns}
\begin{caixaAt}
As cores são somente didáticas. Qualquer atividade em instalação real exige desenergização, bloqueio, documentação e procedimento de segurança aplicável.
\end{caixaAt}
\end{frame}

\begin{frame}{Interruptor simples — contato e bornes explícitos}
\small
\begin{columns}[T,onlytextwidth]
\column{.47\textwidth}
\centering\textbf{Funcional}\par\vspace{1mm}
\begin{circuitikz}[scale=.78,transform shape]
\draw (0,3) node[left]{F} to[spst,l=$S_1$] (2.8,3) to[lamp,l=$L_1$] (2.8,0) -- (0,0) node[left]{N};
\end{circuitikz}
\column{.51\textwidth}
\centering\textbf{Bornes do mecanismo}\par\vspace{1mm}
\begin{tikzpicture}[scale=.78,every node/.style={font=\scriptsize}]
\draw[corNorma,fill=corNorma!5,rounded corners](-2,-1.15)rectangle(2,1.15);
\node at(0,.78){S1 simples};
\node[schterm](c)at(-1.35,-.15){};\node[anchor=north]at(-1.35,-.32){L/C};
\node[schterm](o)at(1.35,-.15){};\node[anchor=north]at(1.35,-.32){1};
\draw[thick](-1.22,-.10)--(1.05,.48);
\draw[LigF,schwire](-3.4,-.15)--(c)node[midway,above]{F};
\draw[LigR,schwire](o)--(3.4,-.15)node[midway,above]{R};
\end{tikzpicture}
\end{columns}
\begin{caixaObs}
O desenho separa o contato interno do percurso dos fios: fase permanente entra em L/C e o borne 1 fornece o retorno para a luminária.
\end{caixaObs}
\end{frame}

\begin{frame}{Ponto simples — fios por trecho físico}
\centering
\begin{tikzpicture}[scale=.81,every node/.style={font=\scriptsize}]
\node[schbox](qd)at(-5.1,1.35){QD\\C1};
\node[schbox,minimum width=2.25cm,minimum height=1cm](ct)at(-1.0,1.35){caixa de teto\\emendas};
\node[schsw](s)at(1.8,-1.4){S1\\L/C \quad 1};
\node[schload](l)at(4.9,1.35){L1};
\draw[very thick,corPrim!20](qd)--(ct);\draw[very thick,corPrim!20](ct)--(s);\draw[very thick,corPrim!20](ct)--(l);
\draw[LigF,schwire]($(qd.east)+(0,.22)$)--($(ct.west)+(0,.22)$);
\draw[LigN,schwire](qd.east)--(ct.west);
\draw[LigPE,schwire]($(qd.east)+(0,-.22)$)--($(ct.west)+(0,-.22)$);
\node[fill=white]at(-3.05,1.88){F + N + PE};
\draw[LigF,schwire](ct.south west)--(-1.0,-1.4)--(s.west)node[pos=.62,below]{F};
\draw[LigR,schwire](s.east)--(3.0,-1.4)--(3.0,.88)--(ct.east)node[pos=.45,right]{R};
\node[fill=white]at(.15,-.83){F + R};
\draw[LigR,schwire](ct.east)--(l.west);
\draw[LigN,schwire](ct.north)--(-1.0,2.5)--(4.9,2.5)--(l.north);
\draw[LigPE,schwire](ct.south east)--(-.15,.35)--(4.9,.35)--(l.south);
\node[fill=white]at(2.35,2.08){R + N + PE};
\end{tikzpicture}
\begin{caixaProjeto}[title=Leitura correta]
O eletroduto até S1 leva fase e retorno; neutro não termina no interruptor simples. O PE acompanha o circuito até as massas previstas, sem ser usado como condutor funcional.
\end{caixaProjeto}
\end{frame}

\begin{frame}{Interruptor duplo — comum e dois retornos independentes}
\centering
\begin{tikzpicture}[scale=.84,every node/.style={font=\scriptsize}]
\draw[corNorma,fill=corNorma!5,rounded corners](-2.2,-1.45)rectangle(2.2,1.45);
\node at(0,1.06){S1 duplo};
\node[schterm](c)at(-1.55,0){};\node[anchor=north]at(-1.55,-.18){L/C};
\node[schterm](o1)at(1.55,.48){};\node[anchor=west]at(1.72,.48){1};
\node[schterm](o2)at(1.55,-.48){};\node[anchor=west]at(1.72,-.48){2};
\draw[thick](-1.42,.04)--(1.28,.50);\draw[thick](-1.42,-.04)--(1.28,-.50);
\draw[LigF,schwire](-4,0)--(c)node[midway,above]{F};
\draw[LigR,schwire](o1)--(4,.48)node[midway,above]{R1};
\draw[LigVum,schwire](o2)--(4,-.48)node[midway,below]{R2};
\node[schload](l1)at(5.2,1.8){L1};\node[schload](l2)at(5.2,-1.8){L2};
\draw[LigR,schwire](4,.48)--(4,1.8)--(l1.west);
\draw[LigVum,schwire](4,-.48)--(4,-1.8)--(l2.west);
\draw[LigN,schwire](-4,2.85)--(5.2,2.85)--(l1.north);\draw[LigN,schwire](5.2,2.85)--(5.2,-1.34);
\draw[LigPE,schwire](-4,-2.85)--(5.65,-2.85)--(5.65,-1.8);\draw[LigPE,schwire](5.65,-2.85)--(5.65,1.8);
\end{tikzpicture}
\begin{caixaAt}
R1 e R2 são saídas diferentes. A ponte de fase entre teclas somente existe quando o conjunto não possui comum interno e a documentação do mecanismo admite essa conexão.
\end{caixaAt}
\end{frame}

\begin{frame}{Interruptores paralelos — contatos SPDT desenhados}
\centering
\begin{tikzpicture}[scale=.79,every node/.style={font=\scriptsize}]
\draw[corNorma,fill=corNorma!5,rounded corners](-5,-1.45)rectangle(-1.8,1.45);
\node at(-3.4,1.08){S1 paralelo};
\node[schterm](c1)at(-4.55,0){};\node[anchor=east]at(-4.75,0){C};
\node[schterm](a1)at(-2.25,.55){};\node[schterm](b1)at(-2.25,-.55){};
\node[anchor=west]at(-2.05,.55){1};\node[anchor=west]at(-2.05,-.55){2};
\draw[thick](-4.42,.04)--(-2.55,.50);
\draw[corNorma,fill=corNorma!5,rounded corners](1.2,-1.45)rectangle(4.4,1.45);
\node at(2.8,1.08){S2 paralelo};
\node[schterm](a2)at(1.65,.55){};\node[schterm](b2)at(1.65,-.55){};\node[schterm](c2)at(3.95,0){};
\node[anchor=east]at(1.45,.55){1};\node[anchor=east]at(1.45,-.55){2};\node[anchor=west]at(4.15,0){C};
\draw[thick](1.95,.50)--(3.82,.04);
\draw[LigF,schwire](-6.8,0)--(c1)node[midway,above]{F};
\draw[LigVum,schwire](a1)--(a2)node[midway,above]{V1};
\draw[LigVdois,schwire](b1)--(b2)node[midway,below]{V2};
\draw[LigR,schwire](c2)--(6.1,0)node[midway,above]{R};
\node[schload](l)at(6.1,2.05){L1};\draw[LigR,schwire](6.1,0)--(l.south);
\draw[LigN,schwire](-6.8,2.95)--(6.1,2.95)--(l.north);
\end{tikzpicture}
\begin{caixaProjeto}[title=O que o aluno precisa enxergar]
F entra no comum de S1; 1 e 2 são os viajantes; o comum de S2 fornece o retorno. A posição da tecla muda qual viajante está ligado ao comum.
\end{caixaProjeto}
\end{frame}

\begin{frame}{Interruptor intermediário — quatro bornes, sem comum}
\centering
\begin{tikzpicture}[scale=.78,every node/.style={font=\scriptsize}]
\node[schsw,minimum width=2cm,minimum height=1.5cm](s1)at(-5,0){S1 paralelo\\C / 1 / 2};
\draw[corNorma,fill=corNorma!5,rounded corners](-1.65,-1.55)rectangle(1.65,1.55);
\node at(0,1.18){S2 intermediário};
\node[schterm](ia1)at(-1.25,.55){};\node[schterm](ia2)at(-1.25,-.55){};\node[schterm](ib1)at(1.25,.55){};\node[schterm](ib2)at(1.25,-.55){};
\node[anchor=east]at(-1.42,.55){1A};\node[anchor=east]at(-1.42,-.55){2A};\node[anchor=west]at(1.42,.55){1B};\node[anchor=west]at(1.42,-.55){2B};
\draw[thick](-1.08,.55)--(1.08,-.55);\draw[thick](-1.08,-.55)--(1.08,.55);
\node[schsw,minimum width=2cm,minimum height=1.5cm](s3)at(5,0){S3 paralelo\\1 / 2 / C};
\draw[LigF,schwire](-6.8,0)--(s1.west)node[midway,above]{F};
\draw[LigVum,schwire]($(s1.east)+(0,.28)$)--(ia1)node[midway,above]{V1A};
\draw[LigVdois,schwire]($(s1.east)+(0,-.28)$)--(ia2)node[midway,below]{V2A};
\draw[LigVum,schwire](ib1)--($(s3.west)+(0,.28)$)node[midway,above]{V1B};
\draw[LigVdois,schwire](ib2)--($(s3.west)+(0,-.28)$)node[midway,below]{V2B};
\draw[LigR,schwire](s3.east)--(6.8,0)node[midway,above]{R};
\end{tikzpicture}
\vspace{2mm}
\begin{caixaObs}
A posição mostrada é cruzada. Na outra posição, 1A liga a 1B e 2A a 2B. O intermediário apenas permuta os viajantes e fica entre dois interruptores paralelos.
\end{caixaObs}
\end{frame}

%=====================================================================
% REPRESENTACAO GRAFICA REVISADA - TOMADAS
%=====================================================================
\section{Leitura gráfica de tomadas}

\begin{frame}{Tomada 2P+T — representação por bornes}
\small
\begin{columns}[T,onlytextwidth]
\column{.52\textwidth}
\centering
\begin{tikzpicture}[scale=.88,every node/.style={font=\scriptsize}]
\draw[corNorma,fill=corNorma!5,rounded corners](-2.2,-1.6)rectangle(2.2,1.6);
\node at(0,1.18){T1 — 2P+T};
\node[circle,draw=corPrim,minimum size=5pt,inner sep=0pt](a)at(-1.05,.2){};
\node[circle,draw=corPrim,minimum size=5pt,inner sep=0pt](b)at(1.05,.2){};
\node[circle,draw=corPrim,minimum size=5pt,inner sep=0pt](c)at(0,-.85){};
\node[anchor=north]at(-1.05,.04){L};
\node[anchor=north]at(1.05,.04){N};
\node[anchor=north]at(0,-1.04){PE};
\draw[thick](-3.4,.2)--(a);
\draw[thick](b)--(3.4,.2);
\draw[thick](0,-2.45)--(c);
\end{tikzpicture}
\column{.46\textwidth}
\begin{caixaProjeto}[title=O que o aluno deve enxergar]
A figura deixa separados os três bornes do mecanismo e evita a ambiguidade criada por um único retângulo representando toda a tomada.
\end{caixaProjeto}
\end{columns}
\begin{caixaObs}
Frente e traseira do mecanismo são vistas diferentes; a identificação impressa no componente deve aparecer no desenho didático.
\end{caixaObs}
\end{frame}

\begin{frame}{Várias tomadas — mostrar cada ponto como derivação}
\centering
\begin{tikzpicture}[scale=.79,every node/.style={font=\scriptsize}]
\node[draw=corPrim,fill=corDef!5,rounded corners,minimum width=1.9cm,minimum height=.85cm](q)at(-5.4,0){origem};
\foreach \x/\n in {-2.6/1,.6/2,3.8/3}{
  \node[draw=corNorma,fill=corNorma!5,rounded corners,minimum width=2.1cm,minimum height=1.15cm](t\n)at(\x,0){T\n\\L \quad N \quad PE};
}
\draw[very thick,corPrim!25](q)--(t1)--(t2)--(t3);
\draw[corPrim,thick]($(q.east)+(0,.28)$)--($(t1.west)+(0,.28)$)--($(t1.east)+(0,.28)$)--($(t2.west)+(0,.28)$)--($(t2.east)+(0,.28)$)--($(t3.west)+(0,.28)$);
\draw[corDef,thick](q.east)--(t1.west)--(t1.east)--(t2.west)--(t2.east)--(t3.west);
\draw[corNorma,thick]($(q.east)+(0,-.28)$)--($(t1.west)+(0,-.28)$)--($(t1.east)+(0,-.28)$)--($(t2.west)+(0,-.28)$)--($(t2.east)+(0,-.28)$)--($(t3.west)+(0,-.28)$);
\end{tikzpicture}
\begin{caixaObs}
O objetivo gráfico é mostrar três caminhos contínuos e três pontos de utilização independentes, em vez de caixas ligadas por uma seta genérica.
\end{caixaObs}
\end{frame}

%=====================================================================
% MONTAGEM PREDIAL - COMPONENTES REAIS
%=====================================================================
\section{Reconhecimento de componentes reais}

\begin{frame}{Ponto de luz — componentes que aparecem na bancada e na obra}
\begin{columns}[T,onlytextwidth]
\column{0.67\textwidth}
\centering
\includegraphics[width=\linewidth,height=.70\textheight,keepaspectratio]{assets/imagens/montagem/componentes_ponto_luz.png}
\column{0.31\textwidth}
\scriptsize
\textbf{Identifique:}
\begin{enumerate}
\item caixa octogonal;
\item soquete/luminária;
\item lâmpada LED;
\item caixa 4$\times$2;
\item interruptor e suporte;
\item placa de acabamento;
\item conectores;
\item condutores por função.
\end{enumerate}
\begin{caixaObs}[title=Regra]
A fotografia reconhece a peça; o esquema define a ligação. A marcação real do borne prevalece sobre a posição aparente.
\end{caixaObs}
\end{columns}
\end{frame}

\begin{frame}{Comandos de iluminação — mecanismos parecidos, funções diferentes}
\begin{columns}[T,onlytextwidth]
\column{0.69\textwidth}
\centering
\includegraphics[width=\linewidth,height=.70\textheight,keepaspectratio]{assets/imagens/montagem/componentes_comandos_iluminacao.png}
\column{0.29\textwidth}
\scriptsize
\textbf{Reconheça:}
\begin{itemize}
\item simples;
\item duplo;
\item paralelo;
\item intermediário;
\item dimmer;
\item sensor;
\item comando eletrônico;
\item bornes traseiros.
\end{itemize}
\begin{caixaAt}
Não identificar um mecanismo apenas pela frente. Conferir símbolo, gravação e diagrama do fabricante.
\end{caixaAt}
\end{columns}
\end{frame}

\begin{frame}{Lâmpadas e drivers LED — o conjunto precisa ser compatível}
\begin{columns}[T,onlytextwidth]
\column{0.69\textwidth}
\centering
\includegraphics[width=\linewidth,height=.70\textheight,keepaspectratio]{assets/imagens/montagem/componentes_led.png}
\column{0.29\textwidth}
\scriptsize
\textbf{Observe:}
\begin{itemize}
\item lâmpadas com driver integrado;
\item luminárias com driver separado;
\item entradas de rede;
\item saídas de baixa tensão;
\item polaridade quando indicada;
\item potência e corrente nominais;
\item condição de dimerização.
\end{itemize}
\begin{caixaObs}
Driver, dimmer e luminária formam um sistema. Compatibilidade elétrica e de controle deve ser verificada em conjunto.
\end{caixaObs}
\end{columns}
\end{frame}

\begin{frame}{Tomada 2P+T — frente, traseira e bornes}
\begin{columns}[T,onlytextwidth]
\column{0.69\textwidth}
\centering
\includegraphics[width=\linewidth,height=.70\textheight,keepaspectratio]{assets/imagens/montagem/componentes_tomada_2pt.png}
\column{0.29\textwidth}
\scriptsize
\textbf{Antes da ligação:}
\begin{itemize}
\item localizar borne de proteção;
\item localizar os dois polos ativos;
\item conferir marcações do fabricante;
\item verificar seção admitida;
\item verificar tipo de conexão;
\item conferir fixação da placa;
\item evitar cobre exposto.
\end{itemize}
\begin{caixaAt}
Vista frontal e vista traseira são opostas. Não deduzir a posição dos bornes olhando apenas a frente do mecanismo.
\end{caixaAt}
\end{columns}
\end{frame}

%=====================================================================
% ESQUEMAS REVISADOS - COMANDOS ESPECIAIS
%=====================================================================
\section{Comandos especiais de iluminação}

\begin{frame}{Relé de impulso — separar circuito de comando e circuito da carga}
\small
\begin{columns}[T,onlytextwidth]
\column{.57\textwidth}
\centering
\begin{tikzpicture}[scale=.76,every node/.style={font=\scriptsize}]
\node[draw=corPrim,fill=corDef!5,rounded corners,minimum width=2cm,minimum height=.8cm](p)at(-3.8,1.55){proteção};
\node[draw=corNorma,fill=corNorma!5,rounded corners,minimum width=2.2cm,minimum height=1.2cm](k)at(0,1.55){contato do relé};
\node[draw=corAccent!85!black,fill=corAccent!7,circle,minimum size=.9cm](l)at(3.8,1.55){L1};
\draw[thick,corPrim](p)--(k)--(l);
\node[draw=corNorma,fill=corNorma!5,rounded corners,minimum width=2.2cm,minimum height=1.2cm](b)at(0,-1.35){bobina\\A1 \quad A2};
\node[draw=corPrim,fill=corDef!5,rounded corners,minimum width=1.5cm,minimum height=.65cm](s1)at(-3.6,-1.35){B1};
\node[draw=corPrim,fill=corDef!5,rounded corners,minimum width=1.5cm,minimum height=.65cm](s2)at(3.6,-1.35){B2};
\draw[thick,corPrim](s1)--(b)--(s2);
\draw[dashed,thick,corAccent,-{Stealth}](b)--(k);
\node[font=\tiny]at(0,.05){acoplamento funcional};
\end{tikzpicture}
\column{.41\textwidth}
\begin{caixaProjeto}[title=Leitura]
As botoeiras pertencem ao comando; o contato do relé pertence ao caminho da carga. O desenho deixa visível que são partes distintas do mesmo dispositivo.
\end{caixaProjeto}
\end{columns}
\end{frame}

\begin{frame}{Sensor de presença — alimentação e saída comutada separadas}
\centering
\begin{tikzpicture}[scale=.84,every node/.style={font=\scriptsize}]
\draw[corNorma,fill=corNorma!5,rounded corners](-1.8,-1.35)rectangle(1.8,1.35);
\node at(0,.88){SENSOR};
\node[circle,draw=corPrim,minimum size=5pt,inner sep=0pt](li)at(-1.2,.15){};
\node[circle,draw=corPrim,minimum size=5pt,inner sep=0pt](ni)at(0,-.55){};
\node[circle,draw=corPrim,minimum size=5pt,inner sep=0pt](lo)at(1.2,.15){};
\node[anchor=north]at(-1.2,0){L};\node[anchor=north]at(0,-.70){N};\node[anchor=north]at(1.2,0){L$'$};
\draw[thick,corPrim](-4,.15)--(li);
\draw[thick,corDef](0,-2.4)--(ni);
\draw[thick,corAccent](lo)--(4,.15);
\node[draw=corAccent!85!black,fill=corAccent!7,circle,minimum size=.9cm](lamp)at(5.0,.15){L1};
\draw[thick,corAccent](4,.15)--(lamp);
\draw[thick,corDef](0,-2.4)--(5,-2.4)--(lamp.south);
\node[font=\tiny,corPrim]at(-3.5,.45){alimentação};
\node[font=\tiny,corAccent]at(3.1,.45){saída comutada};
\end{tikzpicture}
\begin{caixaObs}
A disposição real dos terminais varia; o desenho didático deve separar claramente entrada, referência e saída, em vez de resumir o sensor a uma caixa com uma seta.
\end{caixaObs}
\end{frame}

\begin{frame}{Fotocélula — mesma lógica de leitura por bornes}
\small
\begin{columns}[T,onlytextwidth]
\column{.58\textwidth}
\centering
\begin{tikzpicture}[scale=.78,every node/.style={font=\scriptsize}]
\node[draw=corNorma,fill=corNorma!5,rounded corners,minimum width=2.4cm,minimum height=1.45cm](f)at(0,0){FOTOCÉLULA\\L \quad N \quad L$'$};
\node[draw=corAccent!85!black,fill=corAccent!7,circle,minimum size=.9cm](l)at(4.2,0){L1};
\draw[thick,corPrim](-4,.35)--(f.west |- 0,.35);
\draw[thick,corDef](-4,-.35)--(f.west |- 0,-.35);
\draw[thick,corAccent](f.east)--(l.west);
\draw[thick,corDef](-4,-1.6)--(4.2,-1.6)--(l.south);
\node[font=\tiny]at(-3.4,.65){L};\node[font=\tiny]at(-3.4,-.65){N};\node[font=\tiny]at(2.3,.30){L$'$};
\end{tikzpicture}
\column{.40\textwidth}
\begin{caixaProjeto}[title=Diferença para o sensor]
A função de decisão é luminosa, mas a leitura elétrica continua baseada em alimentação permanente e saída comandada identificada.
\end{caixaProjeto}
\end{columns}
\end{frame}

\begin{frame}{Dimmer e driver LED — mostrar as duas camadas}
\small
\begin{columns}[T,onlytextwidth]
\column{.56\textwidth}
\centering
\begin{tikzpicture}[scale=.74,every node/.style={font=\scriptsize}]
\node[draw=corPrim,fill=corDef!5,rounded corners,minimum width=2.1cm,minimum height=.9cm](rede)at(-4.2,1.1){rede};
\node[draw=corNorma,fill=corNorma!5,rounded corners,minimum width=2.1cm,minimum height=.9cm](drv)at(0,1.1){driver};
\node[draw=corAccent!85!black,fill=corAccent!7,rounded corners,minimum width=2.1cm,minimum height=.9cm](led)at(4.2,1.1){módulo LED};
\draw[thick,corPrim,-{Stealth}](rede)--(drv)--(led);
\node[draw=corNorma,fill=corNorma!5,rounded corners,minimum width=2.1cm,minimum height=.9cm](ctrl)at(0,-1.4){controle};
\draw[dashed,thick,corAccent,<->](ctrl)--(drv);
\node[font=\tiny]at(-2.1,1.45){potência de entrada};
\node[font=\tiny]at(2.1,1.45){saída do driver};
\node[font=\tiny]at(.7,-.25){interface de controle};
\end{tikzpicture}
\column{.42\textwidth}
\begin{itemize}
\item corte de fase;
\item 0--10 V;
\item DALI;
\item PWM/baixa tensão;
\item liga/desliga por relé.
\end{itemize}
\begin{caixaObs}
Um esquema bom não mistura o caminho da potência com a interface de controle. Ambos devem aparecer identificados.
\end{caixaObs}
\end{columns}
\end{frame}

%=====================================================================
% LEITURA E DIAGNOSTICO POR ESQUEMAS
%=====================================================================
\section{Leitura, identificação e diagnóstico}

\begin{frame}{Da prancha à instalação — sequência única de leitura}
\centering
\begin{tikzpicture}[
  etapa/.style={draw=corPrim,fill=corDef!5,rounded corners,minimum width=2.05cm,minimum height=.82cm,align=center,font=\scriptsize},
  fluxo/.style={-{Stealth},thick,corPrim}]
\node[etapa](p)at(-5.2,0){planta\\onde está};
\node[etapa](q)at(-2.6,0){quadro\\de onde vem};
\node[etapa](d)at(0,0){diagrama\\como funciona};
\node[etapa](b)at(2.6,0){bornes\\como conecta};
\node[etapa](e)at(5.2,0){ensaio\\como verifica};
\draw[fluxo](p)--(q)--(d)--(b)--(e);
\end{tikzpicture}
\vspace{4mm}
\begin{caixaProjeto}[title=Critério de interpretação]
O aluno deve conseguir apontar a mesma função em todas as representações: ponto na planta, circuito no quadro, símbolo no diagrama, borne no componente e item do ensaio.
\end{caixaProjeto}
\end{frame}

\begin{frame}{Numeração de fios e bornes — transformar desenho em conexão rastreável}
\small
\begin{columns}[T,onlytextwidth]
\column{.56\textwidth}
\centering
\begin{tikzpicture}[scale=.82,every node/.style={font=\scriptsize}]
\node[draw=corPrim,minimum width=1.8cm](x1)at(-3.1,1.5){X1:01};
\node[draw=corNorma,minimum width=1.8cm](s1)at(0,1.5){S1:L};
\node[draw=corNorma,minimum width=1.8cm](s2)at(0,-.2){S1:1};
\node[draw=corPrim,minimum width=1.8cm](x2)at(3.1,-.2){X2:07};
\draw[thick,corPrim,-{Stealth}](x1)--(s1)node[midway,above]{101};
\draw[thick,corAccent,-{Stealth}](s2)--(x2)node[midway,above]{102};
\draw[dashed,corPrim!60](s1)--(s2);
\node[font=\tiny]at(.55,.65){contato};
\end{tikzpicture}
\column{.42\textwidth}
\begin{itemize}
\item X1/X2: régua ou caixa;
\item 01/07: posição física;
\item S1:L e S1:1: bornes do dispositivo;
\item 101/102: identificação do condutor;
\item a mesma referência deve aparecer na documentação de campo.
\end{itemize}
\end{columns}
\end{frame}

\begin{frame}{Como interpretar um esquema desconhecido}
\small
\begin{enumerate}
\item localizar título, revisão e tensão indicada;
\item marcar origem, proteção e referências de potencial;
\item separar potência, comando, sinal e proteção;
\item seguir um caminho por vez, sem saltar entre ramos;
\item identificar dispositivos pelos seus bornes e contatos;
\item registrar origem e destino de cada fio/cabo;
\item comparar a representação com a marcação do componente;
\item conferir se planta, quadro, diagrama e lista de cabos contam a mesma história.
\end{enumerate}
\begin{caixaAt}
Se um condutor não tiver função e destino claros no desenho, o esquema ainda não está suficientemente compreendido para servir de guia de montagem.
\end{caixaAt}
\end{frame}

\begin{frame}{Diagnóstico por nós — evitar tentativa e erro}
\centering
\begin{tikzpicture}[
  n/.style={draw=corPrim,fill=corDef!5,rounded corners,minimum width=2.15cm,minimum height=.78cm,align=center,font=\scriptsize},
  a/.style={-{Stealth},thick,corPrim}]
\node[n](doc)at(-5.0,1.1){1. conferir\\documentação};
\node[n](vis)at(-2.5,1.1){2. inspeção\\visual};
\node[n](orig)at(0,1.1){3. verificar\\origem};
\node[n](traj)at(2.5,1.1){4. rastrear\\trajeto};
\node[n](carga)at(5.0,1.1){5. verificar\\destino};
\draw[a](doc)--(vis)--(orig)--(traj)--(carga);
\node[n,draw=corNorma,fill=corNorma!5,minimum width=3.3cm](reg)at(0,-1.15){registrar causa, correção\\e nova verificação};
\draw[a](carga.south)--(5.0,-1.15)--(reg.east);
\end{tikzpicture}
\vspace{2mm}
\begin{caixaObs}
O diagrama deve permitir localizar a falha por etapas. Desenhos genéricos demais obrigam o aluno a adivinhar onde começa e termina cada função.
\end{caixaObs}
\end{frame}

%=====================================================================
% APLICACOES ESPECIFICAS - REPRESENTACAO REVISADA
%=====================================================================
\section{Aplicações específicas — leitura por interfaces}

\begin{frame}{Ventilador de teto com luminária — separar duas funções}
\small
\begin{columns}[T,onlytextwidth]
\column{.60\textwidth}
\centering
\begin{tikzpicture}[scale=.75,every node/.style={font=\scriptsize,align=center}]
\node[draw=corPrim,fill=corDef!5,rounded corners,minimum width=2cm,minimum height=.8cm](o)at(-4.1,1.0){origem};
\node[draw=corNorma,fill=corNorma!5,rounded corners,minimum width=2.3cm,minimum height=1.3cm](c)at(-.8,-1.0){controle\\luz \quad ventilador};
\node[draw=corPrim,fill=corDef!5,rounded corners,minimum width=2.2cm,minimum height=1.0cm](t)at(1.6,1.0){caixa de teto};
\node[draw=corAccent!85!black,fill=corAccent!7,circle,minimum size=1.25cm](v)at(5.0,1.0){M +\\L};
\draw[thick,corPrim](o)--(t);
\draw[thick,corPrim](t.south west)--(.4,-1.0)--(c.east);
\draw[thick,corAccent](c.east)+(0,.18)--($(t.south)+(0,-.05)$);
\draw[thick,corNorma](c.east)+(0,-.18)--($(t.south)+(0,.05)$);
\draw[thick,corAccent]($(t.east)+(0,.18)$)--($(v.west)+(0,.18)$);
\draw[thick,corNorma]($(t.east)+(0,-.18)$)--($(v.west)+(0,-.18)$);
\draw[thick,corDef](-4.1,2.25)--(5.0,2.25)--(v.north);
\end{tikzpicture}
\column{.38\textwidth}
\begin{caixaProjeto}[title=Leitura]
O desenho mostra duas saídas de comando chegando ao mesmo conjunto físico. Motor e iluminação não devem aparecer fundidos em uma única seta.
\end{caixaProjeto}
\end{columns}
\end{frame}

\begin{frame}{Exaustor temporizado — distinguir alimentação e gatilho}
\centering
\begin{tikzpicture}[scale=.80,every node/.style={font=\scriptsize,align=center}]
\node[draw=corNorma,fill=corNorma!5,rounded corners,minimum width=2.0cm,minimum height=.85cm](s)at(-4.2,1.35){comando luz};
\node[draw=corAccent!85!black,fill=corAccent!7,circle,minimum size=.9cm](l)at(0,1.35){L1};
\node[draw=corNorma,fill=corNorma!5,rounded corners,minimum width=2.7cm,minimum height=1.3cm](e)at(2.8,-1.05){exaustor temporizado\\L \quad N \quad T};
\draw[thick,corAccent](s)--(l);
\draw[thick,corAccent](s.south)--(-4.2,-1.05)--(e.west)node[midway,below]{gatilho T};
\draw[thick,corPrim](-5.4,-2.2)--(2.8,-2.2)--(e.south west)node[pos=.15,below]{alimentação};
\draw[thick,corDef](-5.4,-2.65)--(3.35,-2.65)--(e.south east)node[pos=.15,below]{referência};
\end{tikzpicture}
\begin{caixaObs}
O temporizador continua sendo uma função própria. A representação revisada separa o sinal de comando da alimentação permanente necessária ao equipamento.
\end{caixaObs}
\end{frame}

\begin{frame}{Ar-condicionado split — força e interligação não são a mesma coisa}
\small
\begin{columns}[T,onlytextwidth]
\column{.60\textwidth}
\centering
\begin{tikzpicture}[scale=.74,every node/.style={font=\scriptsize,align=center}]
\node[draw=corPrim,fill=corDef!5,rounded corners,minimum width=1.9cm,minimum height=.8cm](q)at(-4.5,1.1){origem};
\node[draw=corNorma,fill=corNorma!5,rounded corners,minimum width=2.1cm,minimum height=.8cm](sec)at(-2.0,1.1){seccionamento};
\node[draw=corAccent!85!black,fill=corAccent!7,rounded corners,minimum width=2.2cm,minimum height=1.1cm](ue)at(.8,1.1){unidade\\externa};
\node[draw=corDef,fill=corDef!6,rounded corners,minimum width=2.2cm,minimum height=1.1cm](ui)at(4.3,1.1){unidade\\interna};
\draw[thick,corPrim,-{Stealth}](q)--(sec)--(ue);
\draw[thick,corNorma,<->](ue)--(ui)node[midway,above]{interligação};
\draw[thick,corPrim!45,dashed](-4.5,-.25)--(4.3,-.25);
\node[font=\tiny]at(0,-.55){documentar separadamente origem da força e cabo entre unidades};
\end{tikzpicture}
\column{.38\textwidth}
\begin{caixaProjeto}[title=Representação]
O cabo entre unidades deve aparecer como interface própria, com referência ao diagrama do equipamento, e não como simples prolongamento da alimentação.
\end{caixaProjeto}
\end{columns}
\end{frame}

\begin{frame}{Campainha, interfone e fechadura — separar rede e sistema de sinal}
\centering
\begin{tikzpicture}[scale=.79,every node/.style={font=\scriptsize,align=center}]
\node[draw=corPrim,fill=corDef!5,rounded corners,minimum width=2cm,minimum height=.8cm](r)at(-5.0,1.2){rede};
\node[draw=corNorma,fill=corNorma!5,rounded corners,minimum width=2.2cm,minimum height=.8cm](f)at(-2.2,1.2){fonte/interface};
\node[draw=corDef,fill=corDef!6,rounded corners,minimum width=2.2cm,minimum height=.8cm](p)at(.8,2.0){painel / botoeira};
\node[draw=corDef,fill=corDef!6,rounded corners,minimum width=2.2cm,minimum height=.8cm](i)at(4.2,1.2){interfone};
\node[draw=corDef,fill=corDef!6,rounded corners,minimum width=2.2cm,minimum height=.8cm](fe)at(.8,.1){fechadura};
\draw[thick,corPrim,-{Stealth}](r)--(f);
\draw[thick,corNorma,-{Stealth}](f)--(p)--(i);
\draw[thick,corNorma,-{Stealth}](p)--(fe);
\draw[dashed,very thick,corAccent](-0.7,-1.0)--(-0.7,3.1);
\node[font=\tiny,rotate=90,corAccent]at(-.45,.9){fronteira funcional};
\end{tikzpicture}
\begin{caixaObs}
O esquema deixa visível a fronteira entre alimentação e comunicação/acionamento, evitando representar todo o sistema como um único circuito genérico.
\end{caixaObs}
\end{frame}

\begin{frame}{Iluminação de emergência — alimentação e função de teste separadas}
\centering
\begin{tikzpicture}[scale=.82,every node/.style={font=\scriptsize,align=center}]
\node[draw=corPrim,fill=corDef!5,rounded corners,minimum width=2.1cm,minimum height=.8cm](o)at(-4.2,1.1){alimentação};
\node[draw=corAccent!85!black,fill=corAccent!7,rounded corners,minimum width=2.8cm,minimum height=1.25cm](b)at(0,1.1){bloco autônomo\\rede + bateria};
\node[draw=corNorma,fill=corNorma!5,rounded corners,minimum width=2.1cm,minimum height=.8cm](l)at(4.2,1.1){luz emergência};
\draw[thick,corPrim,-{Stealth}](o)--(b);
\draw[thick,corAccent,-{Stealth}](b)--(l);
\node[draw=corDef,fill=corDef!6,rounded corners,minimum width=2.2cm,minimum height=.75cm](t)at(0,-1.35){teste / supervisão};
\draw[dashed,thick,corNorma,<->](t)--(b);
\end{tikzpicture}
\begin{caixaProjeto}[title=Leitura]
A função de teste/supervisão aparece separada do caminho de energia. Isso facilita relacionar o esquema ao procedimento de verificação do sistema.
\end{caixaProjeto}
\end{frame}

%=====================================================================
% LUMINARIAS E LED - REPRESENTACAO REVISADA
%=====================================================================
\section{Luminárias e sistemas LED}

\begin{frame}{Soquete E27 — representar os dois contatos separadamente}
\small
\begin{columns}[T,onlytextwidth]
\column{.55\textwidth}
\centering
\begin{tikzpicture}[scale=.88,every node/.style={font=\scriptsize}]
\draw[very thick,rounded corners](-2.2,-2.1)rectangle(2.2,2.1);
\draw[corPrim,thick](-1.2,-.6)arc[start angle=200,end angle=-20,radius=1.28];
\draw[corPrim,thick](-1.0,-.1)--(1.0,-.1);
\draw[corPrim,thick](-.82,.36)--(.82,.36);
\draw[corPrim,thick](-.55,.82)--(.55,.82);
\fill[corAccent](0,1.32)circle(.16);
\node[align=center]at(0,1.78){contato\\central};
\node[align=center]at(1.48,.46){contato\\da rosca};
\draw[corAccent,thick](-3.3,1.32)--(-.16,1.32) node[pos=.15,above]{R};
\draw[corDef,thick](3.3,.44)--(1.02,.44) node[pos=.85,above]{N};
\end{tikzpicture}
\column{.43\textwidth}
\begin{caixaProjeto}[title=Por que este desenho é melhor]
Ele mostra que o soquete não é uma carga abstrata com dois fios: existem dois contatos físicos diferentes e o esquema precisa distingui-los.
\end{caixaProjeto}
\end{columns}
\end{frame}

\begin{frame}{Duas luminárias no mesmo comando — dois ramos paralelos}
\centering
\begin{tikzpicture}[scale=.82,every node/.style={font=\scriptsize}]
\node[draw=corNorma,fill=corNorma!5,rounded corners,minimum width=2cm,minimum height=.9cm](s)at(-4.8,0){S1};
\node[draw=corPrim,fill=corDef!5,rounded corners,minimum width=2cm,minimum height=.9cm](cx)at(-1.5,0){derivação};
\node[draw=corAccent!85!black,fill=corAccent!7,circle,minimum size=.9cm](l1)at(2.1,1.45){L1};
\node[draw=corAccent!85!black,fill=corAccent!7,circle,minimum size=.9cm](l2)at(4.9,1.45){L2};
\draw[corPrim,thick](-6.2,0)--(s.west);
\draw[corAccent,thick](s.east)--(cx.west);
\draw[corAccent,thick](cx.east)--(.9,0)--(.9,1.45)--(l1.west);
\draw[corAccent,thick](.9,0)--(3.7,0)--(3.7,1.45)--(l2.west);
\draw[corDef,thick](-6.2,-1.35)--(2.1,-1.35)--(l1.south);
\draw[corDef,thick](2.1,-1.35)--(4.9,-1.35)--(l2.south);
\end{tikzpicture}
\begin{caixaObs}
O desenho mostra uma derivação real em dois ramos. Isso é mais claro do que desenhar duas lâmpadas em sequência na mesma linha.
\end{caixaObs}
\end{frame}

\begin{frame}{Fita LED — fonte, distribuição e trechos devem aparecer}
\small
\begin{columns}[T,onlytextwidth]
\column{.60\textwidth}
\centering
\begin{tikzpicture}[scale=.76,every node/.style={font=\scriptsize}]
\node[draw=corPrim,fill=corDef!5,rounded corners,minimum width=1.8cm,minimum height=.8cm](r)at(-4.8,1.0){rede};
\node[draw=corNorma,fill=corNorma!5,rounded corners,minimum width=2.3cm,minimum height=1.0cm](f)at(-1.6,1.0){fonte / driver};
\node[draw=corPrim,fill=corDef!5,rounded corners,minimum width=2.0cm,minimum height=.9cm](d)at(1.3,1.0){distribuição};
\node[draw=corAccent!85!black,fill=corAccent!7,rounded corners,minimum width=2.4cm,minimum height=.8cm](a)at(4.5,1.75){trecho A};
\node[draw=corAccent!85!black,fill=corAccent!7,rounded corners,minimum width=2.4cm,minimum height=.8cm](b)at(4.5,.25){trecho B};
\draw[thick,corPrim,-{Stealth}](r)--(f)--(d);
\draw[thick,corAccent,-{Stealth}](d)--(a);
\draw[thick,corAccent,-{Stealth}](d)--(b);
\node[font=\tiny]at(-3.2,1.35){entrada};
\node[font=\tiny]at(-.1,1.35){baixa tensão};
\end{tikzpicture}
\column{.38\textwidth}
\begin{caixaProjeto}[title=Leitura]
Trechos longos e derivações não devem ser condensados em uma única linha. O esquema precisa mostrar onde a alimentação é distribuída.
\end{caixaProjeto}
\end{columns}
\end{frame}

\begin{frame}{Driver dimerizável — potência e controle em camadas distintas}
\centering
\begin{tikzpicture}[scale=.80,every node/.style={font=\scriptsize}]
\node[draw=corPrim,fill=corDef!5,rounded corners,minimum width=2cm,minimum height=.85cm](p)at(-4.2,1.1){potência};
\node[draw=corNorma,fill=corNorma!5,rounded corners,minimum width=2.3cm,minimum height=1.1cm](d)at(0,1.1){driver};
\node[draw=corAccent!85!black,fill=corAccent!7,rounded corners,minimum width=2cm,minimum height=.85cm](l)at(4.2,1.1){LED};
\draw[thick,corPrim,-{Stealth}](p)--(d)--(l);
\node[draw=corDef,fill=corDef!6,rounded corners,minimum width=2.3cm,minimum height=.85cm](c)at(0,-1.35){interface de controle};
\draw[dashed,thick,corAccent,<->](c)--(d);
\node[font=\tiny]at(0,-2.0){0--10 V, DALI, PWM ou outra interface prevista};
\end{tikzpicture}
\begin{caixaObs}
Separar graficamente potência e controle facilita entender compatibilidade, endereçamento e diagnóstico do sistema LED.
\end{caixaObs}
\end{frame}

%=====================================================================
% CARGAS DEDICADAS - REPRESENTACAO REVISADA
%=====================================================================
\section{Cargas dedicadas e topologias de alimentação}

\begin{frame}{Mesmo valor de tensão, topologias diferentes}
\small
\begin{columns}[T,onlytextwidth]
\column{.49\textwidth}
\begin{caixaProjeto}[title=Fase--neutro]
\centering
\begin{tikzpicture}[scale=.72,every node/.style={font=\scriptsize,align=center}]
\node[draw=corPrim,fill=corDef!5,rounded corners,minimum width=2cm,minimum height=.8cm](q)at(-2,0){origem};
\node[draw=corAccent!85!black,fill=corAccent!7,rounded corners,minimum width=2cm,minimum height=.8cm](c)at(2,0){carga};
\draw[thick,corPrim]($(q.east)+(0,.18)$)--($(c.west)+(0,.18)$)node[midway,above]{F};
\draw[thick,corDef]($(q.east)+(0,-.18)$)--($(c.west)+(0,-.18)$)node[midway,below]{N};
\end{tikzpicture}
\end{caixaProjeto}
\column{.49\textwidth}
\begin{caixaProjeto}[title=Fase--fase]
\centering
\begin{tikzpicture}[scale=.72,every node/.style={font=\scriptsize,align=center}]
\node[draw=corPrim,fill=corDef!5,rounded corners,minimum width=2cm,minimum height=.8cm](q)at(-2,0){origem};
\node[draw=corAccent!85!black,fill=corAccent!7,rounded corners,minimum width=2cm,minimum height=.8cm](c)at(2,0){carga};
\draw[thick,corPrim]($(q.east)+(0,.18)$)--($(c.west)+(0,.18)$)node[midway,above]{F1};
\draw[thick,corNorma]($(q.east)+(0,-.18)$)--($(c.west)+(0,-.18)$)node[midway,below]{F2};
\end{tikzpicture}
\end{caixaProjeto}
\end{columns}
\begin{caixaObs}
O valor nominal de tensão não informa sozinho quais condutores compõem o circuito. Essa diferença precisa aparecer no diagrama.
\end{caixaObs}
\end{frame}

\begin{frame}{Carga dedicada — mostrar a cadeia física completa}
\centering
\begin{tikzpicture}[scale=.82,every node/.style={font=\scriptsize,align=center}]
\node[draw=corPrim,fill=corDef!5,rounded corners,minimum width=2.1cm,minimum height=.85cm](q)at(-5.0,0){quadro};
\node[draw=corNorma,fill=corNorma!5,rounded corners,minimum width=2.1cm,minimum height=.85cm](p)at(-1.7,0){proteção /\\seccionamento};
\node[draw=corDef,fill=corDef!6,rounded corners,minimum width=2.1cm,minimum height=.85cm](x)at(1.7,0){conexão};
\node[draw=corAccent!85!black,fill=corAccent!7,rounded corners,minimum width=2.1cm,minimum height=.85cm](c)at(5.0,0){carga};
\draw[thick,corPrim,-{Stealth}](q)--(p)--(x)--(c);
\draw[thick,corNorma](-5,-1.35)--(5,-1.35)--(c.south);
\node[font=\tiny]at(0,-1.65){condutor de proteção representado separadamente};
\end{tikzpicture}
\begin{caixaProjeto}[title=Melhoria gráfica]
A representação mostra origem, proteção, ponto de conexão e equipamento como elementos distintos. Isso evita que um único bloco esconda decisões de projeto e interfaces físicas.
\end{caixaProjeto}
\end{frame}

\begin{frame}{Bomba comandada — potência e comando em diagramas coordenados}
\small
\begin{columns}[T,onlytextwidth]
\column{.49\textwidth}
\centering\textbf{Potência}\par\vspace{1mm}
\begin{tikzpicture}[scale=.72,every node/.style={font=\scriptsize,align=center}]
\node[draw=corPrim,fill=corDef!5,minimum width=2cm](q)at(0,2.4){proteção};
\node[draw=corNorma,fill=corNorma!5,minimum width=2cm](k)at(0,1.1){K1};
\node[draw=corAccent!85!black,fill=corAccent!7,circle,minimum size=.95cm](m)at(0,-.5){M};
\draw[thick,corPrim,-{Stealth}](q)--(k)--(m);
\end{tikzpicture}
\column{.49\textwidth}
\centering\textbf{Comando}\par\vspace{1mm}
\begin{tikzpicture}[scale=.72,every node/.style={font=\scriptsize,align=center}]
\node[draw=corPrim,fill=corDef!5,minimum width=2cm](s)at(-2.2,1.3){sensor / boia};
\node[draw=corNorma,fill=corNorma!5,minimum width=2cm](b)at(2.2,1.3){bobina K1};
\draw[thick,corPrim,-{Stealth}](s)--(b);
\draw[dashed,corAccent,-{Stealth}](b.south)--(2.2,-.4)node[below]{atua em K1};
\end{tikzpicture}
\end{columns}
\begin{caixaObs}
A mesma referência K1 liga os dois desenhos. Essa referência cruzada é mais útil ao aluno do que uma única seta entre sensor e motor.
\end{caixaObs}
\end{frame}


%=====================================================================
% MONTAGEM PREDIAL - QUADROS, PROTECOES E BARRAMENTOS
%=====================================================================

\section{Montagem de quadros e dispositivos de proteção}

\begin{frame}{Do unifilar ao quadro físico}
\centering
\begin{tikzpicture}[
  b/.style={draw=corPrim,fill=corDef!7,rounded corners,minimum width=2.0cm,minimum height=.7cm,align=center,font=\scriptsize},
  a/.style={-{Stealth},thick,corPrim}]
\node[b](u)at(-4.5,1.0){Unifilar\\funções e calibres};
\node[b](l)at(-1.5,1.0){Lista de materiais\\modelos e acessórios};
\node[b](a)at(1.5,1.0){Arranjo físico\\módulos e barras};
\node[b](f)at(4.5,1.0){Fiação interna\\bornes e rotas};
\node[b,draw=corNorma,fill=corNorma!7](t)at(0,-1.2){Inspeção, torque, ensaios e identificação};
\draw[a](u)--(l)--(a)--(f);\draw[a](f)--(t);\draw[a](t)--(u);
\end{tikzpicture}
\vspace{4mm}
\begin{caixaAt}
O painel não deve ser “resolvido” durante a montagem. O projeto precisa definir função, corrente, polos, capacidade de interrupção, coordenação, barras, cabos e espaço; a montagem materializa e registra essas decisões.
\end{caixaAt}
\end{frame}

\begin{frame}{Antes de montar: leia a placa e o ambiente}
\scriptsize
\begin{tabularx}{\textwidth}{>{\bfseries}p{2.7cm}X X}
\toprule
Grupo & Conferir & Consequência \\
\midrule
Sistema & tensão, frequência, fases, neutro, PE e esquema de aterramento & polos, barras e ensaios \\
Curto-circuito & corrente presumida no ponto e proteção a montante & Icn/Icu, coordenação e invólucro \\
Corrente/carga & demanda, simultaneidade, harmônicas e reserva & barramentos, cabos e temperatura \\
Ambiente & embutido/sobrepor, umidade, poeira, corrosão e acesso & material, IP, ventilação e posição \\
Operação & responsável, seccionamento e continuidade requerida & agrupamento por DR/RCBO e seletividade \\
Manutenção & expansão, termografia, reaperto e substituição & reserva, acesso e identificação \\
\bottomrule
\end{tabularx}
\begin{caixaAt}
Não instalar quadro cuja corrente nominal ou grau de proteção foram escolhidos apenas pela quantidade de módulos.
\end{caixaAt}
\end{frame}

\begin{frame}{Arquitetura interna de um quadro predial}
\centering
\begin{tikzpicture}[scale=.78,every node/.style={font=\scriptsize}]
\draw[very thick,rounded corners](-6,-3)rectangle(6,3.4);
\draw[fill=black!5](-5.5,2.15)rectangle(5.5,2.75);
\node[draw=corPrim,fill=corDef!7,minimum width=1.6cm,minimum height=.45cm]at(-4.4,2.45){DG};
\node[draw=corAt,fill=corAt!7,minimum width=1.6cm,minimum height=.45cm]at(-2.3,2.45){DPS};
\node[draw=corNorma,fill=corNorma!7,minimum width=1.6cm,minimum height=.45cm]at(-.2,2.45){DR-A};
\node[draw=corNorma,fill=corNorma!7,minimum width=1.6cm,minimum height=.45cm]at(1.9,2.45){DR-B};
\node[draw=corPrim,fill=corDef!7,minimum width=1.6cm,minimum height=.45cm]at(4.0,2.45){reserva};
\draw[very thick](-5.5,1.65)--(5.5,1.65);
\foreach \x/\n in {-4.7/C1,-3.4/C2,-2.1/C3,-.8/C4,.8/C5,2.1/C6,3.4/C7,4.7/C8}{
\node[draw=corPrim,minimum width=1.0cm,minimum height=.7cm]at(\x,.95){\n};}
\draw[very thick](-5.5,.25)--(5.5,.25);
\node[draw=corDef,fill=corDef!7,minimum width=4.4cm,minimum height=.45cm]at(-2.8,-.35){barra N-A};
\node[draw=corDef,fill=corDef!7,minimum width=4.4cm,minimum height=.45cm]at(2.8,-.35){barra N-B};
\node[draw=corNorma,fill=corNorma!7,minimum width=10.0cm,minimum height=.45cm]at(0,-1.35){barra PE};
\node[draw=black!55,fill=black!4,minimum width=10.0cm,minimum height=.45cm]at(0,-2.3){canaleta / terminais / identificação};
\node[font=\tiny,corPrim]at(0,3.05){trilhos DIN e zonas funcionais};
\end{tikzpicture}
\end{frame}

\begin{frame}{Arranjo físico precisa seguir o fluxo elétrico}
\small
\begin{columns}[T,onlytextwidth]
\column{0.48\textwidth}
\begin{caixaProjeto}[title=Fluxo recomendado]
\begin{enumerate}
\item entrada e seccionamento;
\item proteção geral;
\item DPS e sua proteção/coordenação;
\item DR/RCBO ou grupos funcionais;
\item disjuntores terminais;
\item bornes N e PE;
\item saída dos circuitos.
\end{enumerate}
\end{caixaProjeto}
\column{0.50\textwidth}
\begin{itemize}
\item minimizar cruzamentos e laços;
\item separar entrada de saídas quando possível;
\item preservar raio e acesso aos bornes;
\item evitar fontes de calor junto a dispositivos sensíveis;
\item deixar reserva física e elétrica;
\item permitir inspeção sem desmontar metade do quadro;
\item manter diagrama/legenda no invólucro.
\end{itemize}
\end{columns}
\end{frame}

\begin{frame}{Barramento pente — compatibilidade não é só o passo modular}
\small
\begin{columns}[T,onlytextwidth]
\column{0.52\textwidth}
\begin{tikzpicture}[scale=.78]
\foreach \x in {-3,-1.5,0,1.5,3}{\draw[fill=corDef!8,draw=corPrim,thick](\x-.5,0)rectangle(\x+.5,2.3);}
\draw[corAccent,line width=7pt](-3.45,-.45)--(3.45,-.45);
\foreach \x in {-3,-1.5,0,1.5,3}{\draw[corAccent,line width=4pt](\x,-.45)--(\x,0);}
\node[font=\scriptsize]at(0,-1.15){pente de distribuição};
\end{tikzpicture}
\column{0.46\textwidth}
\begin{itemize}
\item corrente nominal e seção do pente;
\item número de fases e sequência dos dentes;
\item distância/passo e formato do borne;
\item linha de produto aceita pelo fabricante;
\item capa e terminação nas pontas livres;
\item alimentação no ponto permitido;
\item torque e inserção total.
\end{itemize}
\end{columns}
\begin{caixaAt}
Não cortar um pente e deixar cobre exposto, nem forçar dente incompatível no borne. Usar acessórios finais e proteção contra toque da mesma solução.
\end{caixaAt}
\end{frame}

\begin{frame}{DR e RCBO — continuidade versus seletividade funcional}
\scriptsize
\begin{columns}[T,onlytextwidth]
\column{0.49\textwidth}
\begin{caixaProjeto}[title=Um DR para vários circuitos]
\begin{tikzpicture}[scale=.62]
\node[draw,minimum width=1.6cm](dr)at(0,2.5){DR-A};
\foreach \x/\n in {-2/C1,0/C2,2/C3}{\node[draw,minimum width=1.3cm](c\n)at(\x,0){\n};\draw[-{Stealth}](dr)--(c\n);}
\end{tikzpicture}
\begin{itemize}
\item menor custo e volume;
\item uma fuga interrompe todo o grupo;
\item soma de fugas naturais precisa ser avaliada.
\end{itemize}
\end{caixaProjeto}
\column{0.49\textwidth}
\begin{caixaProjeto}[title=RCBO por circuito]
\begin{tikzpicture}[scale=.62]
\foreach \x/\n in {-2/C1,0/C2,2/C3}{\node[draw,minimum width=1.3cm,align=center](c\n)at(\x,1.2){RCBO\\\n};\draw[-{Stealth}](c\n)--(\x,-.2);}
\end{tikzpicture}
\begin{itemize}
\item falha fica restrita ao circuito;
\item facilita diagnóstico e continuidade;
\item exige mais módulos e investimento.
\end{itemize}
\end{caixaProjeto}
\end{columns}
\begin{caixaAt}
Tipo, sensibilidade, polos e imunidade do diferencial devem ser definidos pelas cargas e pela norma; não usar “DR genérico” para qualquer circuito eletrônico.
\end{caixaAt}
\end{frame}

\begin{frame}{Cada DR precisa do seu caminho de neutro}
\centering
\begin{tikzpicture}[scale=.82,every node/.style={font=\scriptsize}]
\node[draw=corNorma,fill=corNorma!7,minimum width=1.8cm](a)at(-3.5,2.3){DR-A};
\node[draw=corNorma,fill=corNorma!7,minimum width=1.8cm](b)at(3.5,2.3){DR-B};
\node[draw=corDef,fill=corDef!7,minimum width=3.2cm](na)at(-3.5,0){barra N-A};
\node[draw=corDef,fill=corDef!7,minimum width=3.2cm](nb)at(3.5,0){barra N-B};
\node[draw=corPrim,minimum width=1.5cm](c1)at(-5,-2){C1};
\node[draw=corPrim,minimum width=1.5cm](c2)at(-2,-2){C2};
\node[draw=corPrim,minimum width=1.5cm](c3)at(2,-2){C3};
\node[draw=corPrim,minimum width=1.5cm](c4)at(5,-2){C4};
\draw[corDef,thick,-{Stealth}](a)--(na);\draw[corDef,thick,-{Stealth}](b)--(nb);
\draw[corDef,thick](na)--(c1);\draw[corDef,thick](na)--(c2);\draw[corDef,thick](nb)--(c3);\draw[corDef,thick](nb)--(c4);
\draw[corAt,very thick](-.8,-.1)--(.8,-1.1);\draw[corAt,very thick](-.8,-1.1)--(.8,-.1);
\node[font=\tiny,corAt]at(0,-1.45){não interligar N-A e N-B};
\end{tikzpicture}
\begin{caixaAt}
Fase e neutro de um circuito devem atravessar o mesmo dispositivo diferencial. Neutro compartilhado ou ponte N--PE a jusante causa desarme e elimina a coerência da proteção.
\end{caixaAt}
\end{frame}

\begin{frame}{DPS — conexão curta, coordenada e protegida}
\small
\begin{columns}[T,onlytextwidth]
\column{0.52\textwidth}
\begin{tikzpicture}[scale=.72,every node/.style={font=\scriptsize}]
\node[draw=corPrim,fill=corDef!7,minimum width=2cm](bar)at(-3,1.5){barras F/N};
\node[draw=corAt,fill=corAt!8,minimum width=2cm,minimum height=1.1cm](dps)at(0,1.5){DPS};
\node[draw=corNorma,fill=corNorma!7,minimum width=2cm](pe)at(3,1.5){barra PE};
\draw[corPrim,line width=2pt,-{Stealth}](bar)--(dps);
\draw[corNorma,line width=2pt,-{Stealth}](dps)--(pe);
\draw[dashed,corAccent,<->](-1.9,2.25)--(1.9,2.25);
\node[font=\tiny,corAccent]at(0,2.55){trajeto total mínimo};
\draw[corNorma,thick](pe)--(3,0)node[ground]{};
\end{tikzpicture}
\column{0.46\textwidth}
\begin{itemize}
\item classe/tipo e tensão compatíveis com o sistema;
\item Uc, Up, corrente de descarga e esquema de ligação;
\item proteção de retaguarda/coordenação conforme fabricante;
\item condutores curtos, retos e sem laços;
\item seção e terminal adequados;
\item indicador de fim de vida visível;
\item coordenação com DPS a montante/jusante.
\end{itemize}
\end{columns}
\end{frame}

\begin{frame}{Barra PE — um ponto comum, várias ligações identificadas}
\centering
\begin{tikzpicture}[scale=.82,every node/.style={font=\scriptsize}]
\draw[fill=corNorma!18,draw=corNorma,very thick](-5,-.35)rectangle(5,.35);
\foreach \x in {-4,-2.7,-1.4,0,1.4,2.7,4}{\draw[fill=white,draw=corNorma](\x,0)circle(.14);}
\foreach \x/\n in {-4/C1,-2.7/C2,-1.4/C3,0/invólucro,1.4/DPS,2.7/alimentador,4/reserva}{\draw[corNorma,thick](\x,.35)--(\x,1.55);\node[font=\tiny,rotate=55,anchor=west]at(\x,1.6){\n};}
\draw[corNorma,very thick](2.7,-.35)--(2.7,-1.8);\node[font=\tiny]at(2.7,-2.15){origem/BEP};
\end{tikzpicture}
\vspace{2mm}
\begin{caixaObs}
Cada PE deve ter borne adequado e rastreável. Não empilhar terminais sob um único parafuso sem autorização, não interromper o PE por dispositivo de comando e não usar o trilho DIN como único caminho de proteção.
\end{caixaObs}
\end{frame}

\begin{frame}{Neutro e PE — separados no quadro terminal}
\small
\begin{columns}[T,onlytextwidth]
\column{0.49\textwidth}
\begin{caixaNorma}[title=Barra N]
\begin{itemize}
\item isolada do invólucro quando o esquema exigir;
\item segmentada por DR/RCBO;
\item dimensionada para corrente e harmônicas;
\item cada circuito identificado;
\item sem uso como terra funcional improvisado.
\end{itemize}
\end{caixaNorma}
\column{0.49\textwidth}
\begin{caixaNorma}[title=Barra PE]
\begin{itemize}
\item ligada ao invólucro metálico e ao PE de origem;
\item continuidade permanente;
\item recebe massas, DPS e equipotencialização prevista;
\item sem fusível, chave ou DR em série;
\item conexões acessíveis para inspeção.
\end{itemize}
\end{caixaNorma}
\end{columns}
\begin{caixaAt}
Ponte N--PE em quadro terminal não “melhora o aterramento”: cria corrente no PE, pode impedir o DR de funcionar corretamente e altera o esquema projetado.
\end{caixaAt}
\end{frame}

\begin{frame}{Fiação interna — rota curta sem sacrificar manutenção}
\small
\begin{columns}[T,onlytextwidth]
\column{0.50\textwidth}
\begin{tikzpicture}[scale=.72]
\draw[very thick](-3,-2.7)rectangle(3,2.7);
\draw[fill=black!6](-2.7,1.5)rectangle(2.7,2.0);
\draw[fill=black!6](-2.7,-.3)rectangle(2.7,.2);
\foreach \x in {-2,-1,0,1,2}{\node[draw,minimum width=.7cm,minimum height=.8cm]at(\x,1.75){};\node[draw,minimum width=.7cm,minimum height=.8cm]at(\x,-.05){};}
\draw[corPrim,thick,rounded corners](-2,1.35)--(-2,.75)--(-1,.75)--(-1,.35);
\draw[corFix,thick,rounded corners](0,1.35)--(0,.75)--(1,.75)--(1,.35);
\draw[corNorma,thick,rounded corners](2,-.45)--(2,-1.6)--(-2.5,-1.6);
\node[font=\tiny]at(0,-2.25){canaletas e reserva de curvatura};
\end{tikzpicture}
\column{0.48\textwidth}
\begin{itemize}
\item usar seção e classe compatíveis com os bornes;
\item separar potência, controle e sinais quando aplicável;
\item preservar raio e não tensionar terminais;
\item evitar cabos sobre alavancas e indicadores;
\item não bloquear ventilação ou acesso ao torque;
\item amarrar sem esmagar a isolação;
\item etiquetar nas duas pontas.
\end{itemize}
\end{columns}
\end{frame}

\begin{frame}{Preparação do borne — uma conexão por vez}
\centering
\begin{tikzpicture}[
  b/.style={draw=corPrim,fill=corDef!7,rounded corners,minimum width=2.0cm,minimum height=.7cm,align=center,font=\scriptsize},
  a/.style={-{Stealth},thick,corPrim}]
\node[b](id)at(-5,0){Identificar\\circuito e borne};
\node[b](dec)at(-2.5,0){Decapar\\sem ferir a alma};
\node[b](ter)at(0,0){Aplicar terminal\\se previsto};
\node[b](ins)at(2.5,0){Inserir até\\a referência};
\node[b](tor)at(5,0){Aplicar torque\\e registrar};
\draw[a](id)--(dec)--(ter)--(ins)--(tor);
\end{tikzpicture}
\vspace{5mm}
\begin{caixaAt}
Não dobrar fios para “engrossar”, não estanhar condutor flexível em borne de pressão, não prender isolação e não deixar cobre aparente além do limite do borne.
\end{caixaAt}
\end{frame}

\begin{frame}{Torque — valor, ferramenta e rastreabilidade}
\scriptsize
\begin{tabularx}{\textwidth}{>{\bfseries}p{2.6cm}X X}
\toprule
Etapa & Ação & Registro \\
\midrule
Fonte & obter torque na marcação/ficha do fabricante & referência e revisão consultadas \\
Ferramenta & selecionar torquímetro/chave dentro da faixa & identificação e validade de calibração \\
Aplicação & manter bit correto, alinhamento e um único ciclo controlado & operador, data e posição \\
Marcação & usar selo visual apenas se o procedimento definir & cor/lote e mapa de torque \\
Inspeção & verificar inserção, cobre exposto e deformação & não conformidade e retrabalho \\
Reaperto & somente quando fabricante/procedimento determinar & condição, valor e justificativa \\
\bottomrule
\end{tabularx}
\begin{caixaAt}
Marca de tinta não prova torque. Ela apenas indica que uma etapa foi declarada; a evidência real é ferramenta adequada, valor conhecido e procedimento rastreável.
\end{caixaAt}
\end{frame}

\begin{frame}{Capacidade de interrupção — ler o ponto de instalação}
\small
\begin{columns}[T,onlytextwidth]
\column{0.50\textwidth}
\begin{tikzpicture}[scale=.78,every node/.style={font=\scriptsize}]
\node[draw=corPrim,fill=corDef!7,minimum width=2cm](tr)at(-3,2){fonte/rede};
\node[draw=corPrim,fill=corDef!7,minimum width=2cm](qg)at(0,2){QGBT};
\node[draw=corPrim,fill=corDef!7,minimum width=2cm](qd)at(3,2){QD};
\draw[-{Stealth},thick,corPrim](tr)--(qg)--(qd);
\draw[corAt,thick,-{Stealth}](3,1.35)--(3,-.1);
\node[draw=corAt,fill=corAt!8,minimum width=2.8cm,align=center]at(3,-.7){$I_{cc,pros}$ no QD};
\end{tikzpicture}
\column{0.48\textwidth}
\begin{itemize}
\item obter corrente de curto presumida no barramento;
\item verificar capacidade do conjunto/dispositivo;
\item considerar coordenação ou proteção de retaguarda somente com documentação válida;
\item conferir tensão, polos e categoria do dispositivo;
\item registrar estudo e referência do fabricante;
\item não confundir corrente nominal com capacidade de interrupção.
\end{itemize}
\end{columns}
\end{frame}

\begin{frame}{Balanceamento físico — distribuir e depois medir}
\scriptsize
\begin{columns}[T,onlytextwidth]
\column{0.55\textwidth}
\begin{tabular}{c l r}
\toprule
Fase & Circuitos exemplo & Demanda \\
\midrule
L1 & C1 luz, C4 TUG sala, C7 POS & 4,8 kVA \\
L2 & C2 luz, C5 TUG cozinha, C8 AC & 4,6 kVA \\
L3 & C3 luz, C6 serviço, C9 reserva & 4,7 kVA \\
\bottomrule
\end{tabular}
\column{0.43\textwidth}
\begin{caixaProjeto}[title=Na montagem]
\begin{itemize}
\item seguir sequência do pente/barramento;
\item conferir a fase real de cada módulo;
\item etiquetar no quadro e no circuito;
\item medir correntes em operação representativa;
\item atualizar quadro se houver remanejamento.
\end{itemize}
\end{caixaProjeto}
\end{columns}
\begin{caixaAt}
Uma coluna bem distribuída no papel pode ficar errada se o pente tiver sequência diferente ou se um circuito for movido sem atualizar a fase.
\end{caixaAt}
\end{frame}

\begin{frame}{Lista de materiais do quadro — exemplo rastreável}
\scriptsize
\begin{tabularx}{\textwidth}{c >{\bfseries}p{3.2cm}p{2.1cm}X}
\toprule
Item & Descrição funcional & Quantidade & Dados que fecham a compra \\
\midrule
01 & invólucro/quadro & 1 & módulos úteis, reserva, IP, corrente, porta e acessórios \\
02 & disjuntor geral & 1 & polos, In, curva, Icn/Icu, tensão e linha \\
03 & DPS & conjunto & tipo, Uc, Up, In/Imax, esquema e retaguarda \\
04 & DR/RCBO & por projeto & polos, In, $I_{\Delta n}$, tipo e imunidade \\
05 & disjuntores finais & por circuito & polos, In, curva, capacidade e acessórios \\
06 & barras N/PE & conjunto & corrente, bornes, isolação e segmentação \\
07 & pente/bloco & conjunto & fases, corrente, passo e compatibilidade \\
08 & terminais/etiquetas & lote & seções, classe, impressora e reserva \\
09 & canaletas/trilhos & conjunto & dimensões, tampas e fixação \\
\bottomrule
\end{tabularx}
\begin{caixaObs}
Marca/modelo são fechados depois dos requisitos funcionais, mas a montagem só começa com a compatibilidade real dos componentes confirmada.
\end{caixaObs}
\end{frame}

\begin{frame}{Sequência segura de montagem do quadro}
\scriptsize
\begin{enumerate}
\item receber e conferir invólucro, diagrama, materiais e placas;
\item montar trilhos, canaletas, barras e acessórios mecânicos;
\item posicionar dispositivos conforme arranjo aprovado;
\item validar segregação, acesso, dissipação e reserva;
\item executar barras/pentes e fiação de entrada ainda isolados da fonte;
\item montar saídas, neutros, PE e terminais;
\item identificar todos os dispositivos, cabos e bornes;
\item aplicar torque controlado e inspeção independente;
\item ensaiar continuidade de PE, isolação, polaridade e diferenciais;
\item liberar energização por procedimento e registrar grandezas.
\end{enumerate}
\begin{caixaAt}
Bloqueio e ausência de tensão devem ser confirmados no ponto de trabalho. “Disjuntor geral desligado” sozinho não prova desenergização nem elimina fontes de retorno.
\end{caixaAt}
\end{frame}

\begin{frame}{Inspeção final do quadro — leitura cruzada}
\scriptsize
\begin{tabularx}{\textwidth}{>{\bfseries}p{3.0cm}X X}
\toprule
Documento & Conferir no quadro & Não conformidade típica \\
\midrule
Unifilar & sequência, polos, calibres, DR, DPS e barras & componente diferente sem revisão \\
Quadro de cargas & circuito, fase, seção e uso & etiqueta não corresponde ao ambiente \\
Arranjo & posição, reserva e acessibilidade & módulo invertido/oculto \\
Lista de cabos & origem, destino, formação e PE & cabo ou anilha trocados \\
Mapa de torque & todos os bornes aplicáveis & ponto sem registro \\
Relatório de ensaio & continuidade, isolação, DR e função & resultado sem instrumento/critério \\
As built & alterações autorizadas e data & quadro real diverge da entrega \\
\bottomrule
\end{tabularx}
\begin{caixaProjeto}[title=Resultado esperado]
Outro profissional deve conseguir entender, operar, isolar e manter o quadro sem depender da memória de quem o montou.
\end{caixaProjeto}
\end{frame}

%=====================================================================
% MONTAGEM PREDIAL - ENSAIOS, ENERGIZACAO E COMISSIONAMENTO
%=====================================================================

\section{Ensaios, energização e comissionamento}

\begin{frame}{Comissionar é produzir evidência de que a instalação atende}
\centering
\begin{tikzpicture}[
  b/.style={draw=corPrim,fill=corDef!7,rounded corners,minimum width=2.0cm,minimum height=.7cm,align=center,font=\scriptsize},
  a/.style={-{Stealth},thick,corPrim}]
\node[b](i)at(-5,0){Inspeção\\o que foi montado};
\node[b](e)at(-2.5,0){Ensaios sem tensão\\integridade};
\node[b](en)at(0,0){Energização\\controlada};
\node[b](f)at(2.5,0){Ensaios funcionais\\desempenho};
\node[b](r)at(5,0){Registros\\e entrega};
\draw[a](i)--(e)--(en)--(f)--(r);
\end{tikzpicture}
\vspace{5mm}
\begin{caixaAt}
“Funcionou” não é evidência suficiente. Cada resultado precisa de ponto ensaiado, método, instrumento, valor, critério, responsável, data e vínculo com a revisão do projeto.
\end{caixaAt}
\end{frame}

\begin{frame}{Plano de inspeção e ensaios — ITP}
\scriptsize
\begin{tabularx}{\textwidth}{c >{\bfseries}p{2.7cm}X X}
\toprule
Etapa & Verificação & Critério de aceitação & Registro \\
\midrule
Recebimento & cabos, quadros, dispositivos e certificados & especificação/revisão aprovada & ficha e rastreio de lote \\
Infraestrutura & rota, caixas, segregação e acesso & projeto e inspeção antes do fechamento & checklist/fotos \\
Cabos & seção, identificação, terminação e torque & projeto + fabricante & lista de cabos/mapa \\
Sem tensão & continuidade, isolação e polaridade & procedimento/norma aplicável & folha de medição \\
Com tensão & tensão, sequência, DR, função e carga & projeto, norma e fabricante & relatório de ensaio \\
Entrega & as built, ajustes e pendências & dossiê completo e aprovado & termo de comissionamento \\
\bottomrule
\end{tabularx}
\begin{caixaObs}
Pontos de espera impedem fechar parede, forro ou quadro antes da inspeção que depois ficaria impossível.
\end{caixaObs}
\end{frame}

\begin{frame}{Instrumento certo, faixa certa e estado conhecido}
\scriptsize
\begin{tabularx}{\textwidth}{>{\bfseries}p{2.6cm}X X}
\toprule
Instrumento & Uso & Conferir antes do ensaio \\
\midrule
Detector de tensão & constatação conforme procedimento & categoria, faixa e prova antes/depois \\
Multímetro & tensão, resistência e continuidade geral & CAT, pontas, fusíveis e função selecionada \\
Micro-ohmímetro/baixa R & continuidade de PE e ligações & corrente de teste, zero dos cabos e resolução \\
Megômetro & resistência de isolação & tensão de ensaio e equipamentos desconectados \\
Testador de DR & corrente/tempo de atuação & tipo de DR, sequência e condição do circuito \\
Medidor de loop & impedância e curto presumido & método compatível com DR e sistema \\
Terrômetro & resistência de eletrodo & método, geometria e interferências \\
Luxímetro/alicate & desempenho e corrente & faixa, orientação, calibração e posição \\
\bottomrule
\end{tabularx}
\begin{caixaAt}
Categoria CAT não substitui técnica de trabalho. Instrumento inadequado, cabo danificado ou função errada pode criar arco, choque e leitura falsa.
\end{caixaAt}
\end{frame}

\begin{frame}{Sequência antes de qualquer medição energizada}
\scriptsize
\begin{enumerate}
\item revisar análise de risco, procedimento e limites do trabalho;
\item identificar todas as fontes, inclusive gerador, UPS, fotovoltaico e retornos;
\item seccionar, impedir reenergização e sinalizar;
\item constatar ausência de tensão com método/instrumento previstos;
\item descarregar/aterrar quando aplicável;
\item delimitar a área e proteger partes vizinhas;
\item executar inspeção e ensaios desenergizados;
\item retirar curtos, pontes e instrumentos de ensaio;
\item recompor barreiras e confirmar prontidão;
\item autorizar energização controlada com equipe informada.
\end{enumerate}
\begin{caixaAt}
Esta sequência não substitui o procedimento da organização nem a NR-10. O responsável define o método específico para a instalação e a equipe autorizada.
\end{caixaAt}
\end{frame}

\begin{frame}{Inspeção visual — defeitos que o instrumento pode não revelar}
\small
\begin{columns}[T,onlytextwidth]
\column{0.49\textwidth}
\begin{itemize}
\item componentes e seções conforme projeto;
\item proteção contra contato direto e barreiras;
\item identificação de circuitos, neutro e PE;
\item separação, IP e integridade de caixas;
\item conexão, terminal e torque registrados;
\item barras N por DR e PE contínuo;
\item dispositivos acessíveis e sem obstrução;
\item materiais sem dano, umidade ou contaminação.
\end{itemize}
\column{0.49\textwidth}
\begin{caixaProjeto}[title=Leitura cruzada]
Comparar no local:
\begin{enumerate}
\item planta e circuito;
\item quadro de cargas;
\item unifilar;
\item lista de cabos;
\item detalhe de montagem;
\item revisão liberada;
\item alteração/as built.
\end{enumerate}
\end{caixaProjeto}
\end{columns}
\end{frame}

\begin{frame}{Continuidade do PE — medir o caminho completo}
\scriptsize
\begin{columns}[T,onlytextwidth]
\column{0.52\textwidth}
\begin{tikzpicture}[scale=.76,every node/.style={font=\scriptsize}]
\node[draw=corNorma,fill=corNorma!7,minimum width=2cm](bar)at(-3.6,1.5){barra PE};
\node[draw=corPrim,minimum width=2cm](massa)at(3.6,1.5){massa remota};
\node[draw=corAccent,fill=corAccent!7,minimum width=2.2cm](m)at(0,-1.0){medidor baixa R};
\draw[corNorma,line width=2.2pt](bar)--(massa);
\draw[corAccent,thick](m.west)--(-3.6,-1)--(bar.south);
\draw[corAccent,thick](m.east)--(3.6,-1)--(massa.south);
\node[font=\tiny,corNorma]at(0,1.82){PE do circuito};
\end{tikzpicture}
\column{0.46\textwidth}
\begin{enumerate}
\item zerar/subtrair resistência dos cabos de teste;
\item medir do barramento PE a cada massa e tomada;
\item registrar comprimento/seção quando relevante;
\item comparar com resistência esperada e critérios do projeto;
\item investigar valores instáveis ou crescentes;
\item repetir após qualquer retrabalho.
\end{enumerate}
\end{columns}
\begin{caixaAt}
O bip de continuidade de um multímetro pode aceitar dezenas de ohms. Para provar conexão protetiva, usar método e resolução adequados ao baixo valor esperado.
\end{caixaAt}
\end{frame}

\begin{frame}{Resistência de isolação — preparar o circuito antes do megômetro}
\scriptsize
\centering
\begin{tikzpicture}[scale=.72,every node/.style={font=\scriptsize}]
\node[draw=corPrim,fill=corDef!7,minimum width=2cm](meg)at(-4.5,0){megômetro};
\node[draw=corPrim,minimum width=2cm,align=center](ativos)at(0,1.2){ativos unidos\\F + N};
\node[draw=corNorma,minimum width=2cm](pe)at(0,-1.2){PE / massa};
\node[draw=corAt,fill=corAt!7,minimum width=2.5cm,align=center](eq)at(4.5,0){eletrônica e DPS\\isolados conforme método};
\draw[corPrim,thick,-{Stealth}](meg)--(ativos);
\draw[corNorma,thick,-{Stealth}](meg)--(pe);
\draw[corAt,dashed,very thick](2.4,-1.8)rectangle(5.8,1.8);
\draw[corAt,very thick](2.8,-1.3)--(5.4,1.3);\draw[corAt,very thick](2.8,1.3)--(5.4,-1.3);
\end{tikzpicture}
\vspace{0.5mm}
\begin{enumerate}
\item confirmar ausência de tensão e descarregar capacitores;
\item desconectar/proteger equipamentos que não suportam a tensão de ensaio;
\item definir grupos de condutores e tensão de teste no procedimento aplicável;
\item medir, aguardar estabilização e registrar valor/tempo/temperatura;
\item descarregar o circuito após o ensaio e recompor conexões.
\end{enumerate}
\begin{caixaAt}
Aplicar 500 Vcc indiscriminadamente pode danificar driver, SPD, automação e equipamentos eletrônicos. O plano de ensaio precisa dizer o que fica conectado e qual tensão usar.
\end{caixaAt}
\end{frame}

\begin{frame}{Critérios de isolação — exemplo de tabela de procedimento}
\scriptsize
\begin{tabularx}{\textwidth}{>{\bfseries}p{3.1cm}c c X}
\toprule
Circuito ensaiado & Tensão cc de referência & Limite de referência & Condição obrigatória \\
\midrule
SELV/PELV & 250 V & 0,5 M$\Omega$ & fontes/equipamentos tratados conforme o fabricante \\
Até 500 V, inclusive FELV & 500 V & 1,0 M$\Omega$ & eletrônica e DPS isolados quando necessário \\
Acima de 500 V & 1 000 V & 1,0 M$\Omega$ & procedimento específico e componentes compatíveis \\
\bottomrule
\end{tabularx}
\begin{caixaNorma}[title=Como usar esta tabela]
Valores são referência didática tradicional de verificação de BT. Antes do ensaio real, confirmar a edição vigente da norma aplicável, o procedimento da organização e as limitações de todos os equipamentos conectados.
\end{caixaNorma}
\begin{caixaAt}
Resultado alto não corrige circuito ligado no borne errado. Isolação e polaridade são verificações diferentes.
\end{caixaAt}
\end{frame}

\begin{frame}{Polaridade — o comando deve interromper o condutor previsto}
\small
\begin{columns}[T,onlytextwidth]
\column{0.50\textwidth}
\begin{tikzpicture}[scale=.75]
\draw[corPrim,thick](-3,2)--(-1.2,2);\node[font=\tiny]at(-2.6,2.3){F};
\node[draw=corNorma,minimum width=1.8cm](s)at(0,2){S1};
\node[draw,circle,minimum size=.9cm](l)at(3,2){L1};
\draw[corAccent,thick](s)--(l);
\draw[corDef,thick](-3,0)--(3,0)--(l.south);\node[font=\tiny]at(-2.6,.3){N};
\draw[corNorma,thick](-3,-.7)--(3.5,-.7)--(3.5,2);\node[font=\tiny]at(-2.5,-1){PE};
\draw[corNorma,thick](-1.2,2)--(s);
\end{tikzpicture}
\column{0.48\textwidth}
Verificar, conforme o sistema:
\begin{itemize}
\item fase no borne previsto de interruptores e proteções;
\item polos ativos seccionados juntos quando exigido;
\item tomadas e conectores segundo marcação/padrão do projeto;
\item neutro no barramento do DR correto;
\item PE no contato de proteção e massas;
\item ausência de pontes N--PE indevidas.
\end{itemize}
\end{columns}
\end{frame}

\begin{frame}{Desligamento automático — relacionar loop e proteção}
\small
Para um esquema em que o critério seja aplicável:
\[
Z_s\leq\frac{U_0}{I_a}
\]
\begin{columns}[T,onlytextwidth]
\column{0.48\textwidth}
\begin{caixaDef}[title=Termos]
\begin{itemize}
\item $Z_s$: impedância do laço de falta;
\item $U_0$: tensão fase--terra;
\item $I_a$: corrente que garante atuação no tempo requerido.
\end{itemize}
\end{caixaDef}
\column{0.50\textwidth}
\begin{caixaEx}[title=Exemplo didático]
Para $U_0=220$ V e $I_a=100$ A:
\[
Z_{s,max}=\frac{220}{100}=2{,}20\ \Omega.
\]
Se a medição corrigida for 1,10 $\Omega$, o critério numérico é atendido.
\end{caixaEx}
\end{columns}
\begin{caixaAt}
$I_a$ não é simplesmente a corrente nominal do disjuntor. Deve vir da característica tempo--corrente e do tempo de desligamento exigido; DR pode participar da medida de proteção conforme o esquema.
\end{caixaAt}
\end{frame}

\begin{frame}{Teste de DR — botão T e instrumento respondem perguntas diferentes}
\scriptsize
\begin{tabularx}{\textwidth}{>{\bfseries}p{2.9cm}X X}
\toprule
Verificação & O que demonstra & Limitação \\
\midrule
Botão T do DR & mecanismo e circuito interno de teste têm condição de atuar & não mede tempo, corrente nem PE da instalação \\
Testador no ponto & atuação para corrente/fase programadas e tempo medido & exige condição correta do circuito e método seguro \\
Inspeção de tipo & DR correto para a forma de onda/carga & depende da especificação e marcação \\
Rastreio de N & neutro passa pelo mesmo DR e não é compartilhado & precisa inspeção/continuidade do circuito \\
Fuga operacional & margem entre fuga natural e sensibilidade & varia com equipamentos e regime \\
\bottomrule
\end{tabularx}
\begin{caixaAt}
Não repetir disparos indefinidamente sem investigar. Registrar corrente de teste, ângulo/fase quando aplicável, tempo, ponto ensaiado e critério da edição normativa/procedimento.
\end{caixaAt}
\end{frame}

\begin{frame}{Aterramento — o método depende do que se quer medir}
\small
\begin{columns}[T,onlytextwidth]
\column{0.51\textwidth}
\begin{tikzpicture}[scale=.72,every node/.style={font=\scriptsize}]
\draw[very thick](-5,0)--(5,0);
\draw[corNorma,line width=3pt](-3,0)--(-3,-2.2);\node[font=\tiny]at(-3,-2.5){eletrodo E};
\draw[corAccent,line width=2pt](0,0)--(0,-1.4);\node[font=\tiny]at(0,-1.7){sonda P};
\draw[corPrim,line width=2pt](3.5,0)--(3.5,-1.4);\node[font=\tiny]at(3.5,-1.7){sonda C};
\node[draw=corDef,fill=corDef!7,minimum width=2cm](t)at(-.2,2){terrômetro};
\draw[corNorma,thick](t)--(-3,0);\draw[corAccent,thick](t)--(0,0);\draw[corPrim,thick](t)--(3.5,0);
\end{tikzpicture}
\column{0.47\textwidth}
\begin{itemize}
\item queda de potencial mede o eletrodo com geometria adequada;
\item alicate mede laço existente, não um eletrodo isolado qualquer;
\item continuidade prova a ligação PE, não a resistência do solo;
\item impedância de loop avalia caminho de falta do sistema;
\item registrar umidade, posição das sondas e conexões existentes;
\item comparar ao critério do esquema de aterramento/projeto.
\end{itemize}
\end{columns}
\end{frame}

\begin{frame}{Energização controlada — etapas e pontos de parada}
\centering
\begin{tikzpicture}[
  b/.style={draw=corPrim,fill=corDef!7,rounded corners,minimum width=1.8cm,minimum height=.68cm,align=center,font=\tiny},
  a/.style={-{Stealth},thick,corPrim}]
\node[b](r)at(-5.1,0){Revisão final\\barreiras e cabos};
\node[b](g)at(-3.05,0){Energizar DG\\sem cargas};
\node[b](v)at(-1.0,0){Medir tensões\\e sequência};
\node[b](c)at(1.05,0){Energizar\\circuito a circuito};
\node[b](f)at(3.1,0){Testar função\\e proteção};
\node[b](m)at(5.15,0){Aplicar carga\\e monitorar};
\draw[a](r)--(g)--(v)--(c)--(f)--(m);
\end{tikzpicture}
\vspace{5mm}
\begin{caixaAt}[title=Parar imediatamente se houver]
Tensão inesperada, ruído, cheiro, aquecimento, desarme sem causa conhecida, corrente incompatível, indicação errada, perda de fase/neutro ou qualquer risco não controlado. Desenergizar pelo procedimento e investigar antes de repetir.
\end{caixaAt}
\end{frame}

\begin{frame}{Tensão, sequência de fases e rotação}
\scriptsize
\begin{tabularx}{\textwidth}{>{\bfseries}p{2.7cm}X X}
\toprule
Medição & Objetivo & Erro evitado \\
\midrule
F--N em cada fase & confirmar tensão de utilização e neutro & equipamento 127/220/230 V em tensão errada \\
F--F entre pares & confirmar sistema trifásico e simetria & fase ausente ou ligação incorreta \\
N--PE & detectar diferença anormal sob condição conhecida & neutro frouxo/corrente indevida no PE \\
Sequência L1-L2-L3 & confirmar ordem de fases & inversão de motores/elevadores \\
Corrente por fase/N & conferir carga e equilíbrio & sobrecarga ou harmônicas no neutro \\
\bottomrule
\end{tabularx}
\begin{caixaAt}
Corrigir rotação no ponto previsto do circuito, com instalação desenergizada. Não trocar fases aleatoriamente em vários quadros e perder a sequência documentada.
\end{caixaAt}
\end{frame}

\begin{frame}{Ensaio funcional — verificar comando, falha e retorno}
\scriptsize
Para cada sistema, registrar:
\begin{enumerate}
\item condição inicial e permissivos;
\item comando local, remoto e automático;
\item indicação de ligado/desligado/falha;
\item comportamento na perda e no retorno de energia;
\item atuação de intertravamentos e limites;
\item parada normal e de emergência quando aplicável;
\item modo manual e automático sem comandos simultâneos;
\item ajuste final de temporizadores, sensores e cenas;
\item corrente, ruído, temperatura e anomalias;
\item coerência entre função real e diagrama.
\end{enumerate}
\begin{caixaObs}
Ensaiar também a condição de falha prevista. Um sistema que só funciona no cenário ideal ainda não foi comissionado.
\end{caixaObs}
\end{frame}

\begin{frame}{Carga e termografia — resultado depende da condição}
\small
\begin{columns}[T,onlytextwidth]
\column{0.48\textwidth}
\begin{caixaNorma}[title=Registrar a condição]
\begin{itemize}
\item corrente e percentual de carga;
\item tempo desde a energização;
\item temperatura ambiente e ventilação;
\item emissividade/distância/reflexos;
\item ponto comparável entre fases;
\item imagem térmica + imagem visível;
\item equipamento e ajuste usados.
\end{itemize}
\end{caixaNorma}
\column{0.50\textwidth}
\begin{caixaAt}[title=Não concluir pelo “ponto vermelho”]
Temperatura aparente isolada não prova defeito. Avaliar diferença entre componentes equivalentes, carga, limite do fabricante, conexão, ventilação e tendência. Corrigir a causa e repetir sob condição comparável.
\end{caixaAt}
\end{columns}
\end{frame}

\begin{frame}{Folha de resultados — exemplo mínimo}
\scriptsize
\begin{tabularx}{\textwidth}{p{1.2cm}p{2.1cm}p{1.8cm}p{1.8cm}p{1.8cm}X}
\toprule
ID & Ponto/circuito & Ensaio & Resultado & Critério & Situação/ação \\
\midrule
E-01 & QD1--T03 & PE baixa R & 0,18 $\Omega$ & procedimento P-04 & conforme \\
E-02 & C7 & isolação & 250 M$\Omega$ & procedimento P-05 & conforme \\
E-03 & T03 & polaridade & L/N/PE corretos & diagrama R03 & conforme \\
E-04 & C7 & DR & 22 ms & critério vigente & conforme \\
E-05 & QD1 & tensão L1-N & 221,4 V & faixa de projeto & conforme \\
E-06 & L12 & função & falha intermitente & sequência F-12 & refazer borne X2:07 \\
\bottomrule
\end{tabularx}
\vspace{2mm}
\begin{caixaObs}
Anexar instrumento, número de série, calibração, responsável, data, revisão do projeto, temperatura/condição e evidência do retrabalho com novo ensaio.
\end{caixaObs}
\end{frame}

\begin{frame}{Diagnóstico: do sintoma à causa comprovada}
\centering
\begin{tikzpicture}[scale=0.78,transform shape,
  b/.style={draw=corPrim,fill=corDef!7,rounded corners,minimum width=2.05cm,minimum height=.72cm,align=center,font=\scriptsize},
  a/.style={-{Stealth},thick,corPrim}]
\node[b](s)at(0,3.2){Definir sintoma\\e condição};
\node[b](seg)at(-3.5,1.7){Tornar seguro\\e preservar evidência};
\node[b](doc)at(3.5,1.7){Ler diagrama\\e histórico};
\node[b](hip)at(0,.1){Criar hipóteses\\ordenadas};
\node[b](test)at(-3.5,-1.5){Testar uma\\hipótese por vez};
\node[b](corr)at(0,-2.8){Corrigir causa\\e atualizar documento};
\node[b](re)at(3.5,-1.5){Reensaiar\\função e proteção};
\draw[a](s)--(seg);\draw[a](s)--(doc);\draw[a](seg)--(hip);\draw[a](doc)--(hip);\draw[a](hip)--(test)--(corr)--(re);
\end{tikzpicture}
\begin{caixaAt}
Troca aleatória de componentes cria novas falhas e apaga evidências. A causa só está comprovada quando explica o sintoma e o reensaio fecha sem efeitos colaterais.
\end{caixaAt}
\end{frame}

\begin{frame}{Dossiê de entrega e operação}
\small
\begin{columns}[T,onlytextwidth]
\column{0.49\textwidth}
\begin{itemize}
\item projeto as built e revisões encerradas;
\item memoriais e quadros de cargas finais;
\item diagramas, listas de cabos e bornes;
\item certificados/fichas dos componentes;
\item mapas de torque e inspeção;
\item relatórios de continuidade, isolação, DR, aterramento e função;
\item pendências e desvios formalmente aceitos;
\item treinamento e contatos de suporte.
\end{itemize}
\column{0.49\textwidth}
\begin{caixaProjeto}[title=Plano de manutenção]
\begin{itemize}
\item periodicidade por criticidade/ambiente;
\item inspeções e testes recorrentes;
\item limpeza e acesso seguro;
\item critérios de substituição;
\item peças sobressalentes;
\item controle de revisão;
\item registro de falhas e intervenções.
\end{itemize}
\end{caixaProjeto}
\end{columns}
\end{frame}

\begin{frame}{Avaliação prática — comissionar um circuito completo}
\scriptsize
O aluno recebe planta, unifilar, quadro de cargas, esquema de ligação e circuito montado com uma falha. Deve:
\begin{enumerate}
\item validar revisão, fontes e análise de risco;
\item inspecionar a montagem contra os documentos;
\item propor o plano e os critérios de ensaio;
\item provar ausência de tensão conforme procedimento;
\item medir PE, isolação e polaridade;
\item localizar e corrigir a falha com autorização;
\item recompor, energizar de forma controlada e testar função;
\item preencher relatório e atualizar o as built;
\item explicar quais evidências permitem liberar o circuito.
\end{enumerate}
\begin{caixaAt}
Critério de aprovação: nenhuma etapa de segurança pode ser compensada por um resultado funcional correto.
\end{caixaAt}
\end{frame}

%=====================================================================
% MONTAGEM PREDIAL - PRATICA INTEGRADA DO PROJETO A ENTREGA
%=====================================================================

\section{Prática integrada — do projeto à execução}

\newcommand{\UnidadeDidaticaBase}{%
\draw[very thick](0,0)rectangle(8,6);
\draw[thick](0,2.4)--(8,2.4);
\draw[thick](5.0,2.4)--(5.0,6);
\draw[thick](2.8,0)--(2.8,2.4);
\draw[white,line width=5pt](1.0,0)--(2.0,0);\draw(1.0,0)--(2.0,1.0);
\draw[white,line width=5pt](1.0,2.4)--(1.9,2.4);\draw(1.0,2.4)--(1.9,3.3);
\draw[white,line width=5pt](5.5,2.4)--(6.4,2.4);\draw(5.5,2.4)--(6.4,3.3);
\draw[white,line width=5pt](2.8,.7)--(2.8,1.55);\draw(2.8,.7)--(3.65,1.55);
\node[font=\tiny\bfseries]at(2.5,5.5){SALA / QUARTO};
\node[font=\tiny\bfseries]at(6.5,5.5){COPA};
\node[font=\tiny\bfseries]at(1.4,1.95){BANHEIRO};
\node[font=\tiny\bfseries]at(5.4,1.2){ENTRADA / CIRCULAÇÃO};}

\begin{frame}{Desafio integrado — quatro competências no mesmo trabalho}
\centering
\begin{tikzpicture}[
  b/.style={draw=corPrim,fill=corDef!7,rounded corners,minimum width=2.25cm,minimum height=.78cm,align=center,font=\scriptsize},
  a/.style={-{Stealth},thick,corPrim}]
\node[b](l)at(-4.2,1){Ler e interpretar\\documentos};
\node[b](p)at(-1.4,-.7){Projetar\\e dimensionar};
\node[b](e)at(1.4,1){Executar\\e identificar};
\node[b](c)at(4.2,-.7){Ensaiar\\e entregar};
\draw[a](l)--(p)--(e)--(c);
\draw[a](c)to[bend left=30](l);
\end{tikzpicture}
\vspace{4mm}
\begin{caixaProjeto}[title=Regra da prática]
Toda alteração descoberta na execução volta ao projeto; todo circuito projetado termina em inspeção, ensaio e registro. Não existem notas separadas para desenho “bonito” e montagem que funciona por acaso.
\end{caixaProjeto}
\end{frame}

\begin{frame}{Unidade didática — arquitetura e programa}
\centering
\begin{tikzpicture}[scale=.82,every node/.style={font=\scriptsize}]
\UnidadeDidaticaBase
\draw[fill=black!5](.45,3.1)rectangle(2.9,4.0);\node[font=\tiny]at(1.68,3.55){cama/sofá};
\draw[fill=black!5](3.7,4.4)rectangle(4.6,5.7);\node[font=\tiny,rotate=90]at(4.15,5.05){armário};
\draw[fill=black!5](5.3,4.45)rectangle(7.7,5.05);\node[font=\tiny]at(6.5,4.75){bancada};
\draw[fill=black!5](5.4,3.0)rectangle(6.3,3.9);\node[font=\tiny]at(5.85,3.45){gel.};
\draw[fill=black!5](.4,.35)rectangle(1.5,1.2);\node[font=\tiny]at(.95,.78){box};
\end{tikzpicture}
\vspace{1mm}
\begin{caixaObs}
Módulo de 48 m$^2$: iluminação com comando simples, paralelo e sensor; TUG gerais e de bancada; aquecedor 5,5 kW; ar-condicionado 1,2 kW; quadro próximo ao acesso.
\end{caixaObs}
\end{frame}

\begin{frame}{Planta elétrica — ler localização, circuito e comando}
\centering
\begin{tikzpicture}[scale=.92,every node/.style={font=\tiny}]
\UnidadeDidaticaBase
\SimQD{7.55}{1.25}{QD1}
\SimLuz{2.5}{4.4}\SimLuz{6.5}{4.4}\SimLuz{1.4}{1.2}\SimLuz{5.3}{1.2}
\SimInterruptor{1.1}{2.75}{S1p}\SimInterruptor{4.65}{2.75}{S2p}\SimInterruptor{5.4}{2.75}{S3}\SimInterruptor{2.45}{.45}{S4}
\foreach \x/\y in {.25/3.1,.25/5.4,4.75/3.1,4.75/5.4,5.25/5.8,7.75/5.8}{\SimTUG{\x}{\y}}
\foreach \x in {5.5,6.5,7.5}{\SimTUGM{\x}{5.1}}
\SimTUE{1.45}{.35}{AQ}\SimTUE{4.65}{5.65}{AC}
\draw[corAccent,dashed,thick](7.25,1.25)--(5.3,1.2)--(1.4,1.2)--(2.5,4.4)--(6.5,4.4);
\draw[corPrim,dashed,thick](7.25,1.0)--(7.2,3.1)--(4.75,3.1)--(.25,3.1);
\draw[corFix,dashed,thick](7.25,.75)--(1.45,.35);
\node[fill=white,corAccent,font=\tiny\bfseries]at(4.8,1.5){C1};
\node[fill=white,corPrim,font=\tiny\bfseries]at(6.8,3.4){C2/C3};
\node[fill=white,corFix,font=\tiny\bfseries]at(4,.6){C4};
\end{tikzpicture}
\vspace{1mm}
\scriptsize S1p/S2p: comando paralelo da sala; S3: copa; S4: banheiro; circulação por sensor no ponto L4. Rotas são didáticas e devem ser convertidas em comprimentos reais.
\end{frame}

\begin{frame}{Quadro de cargas — ponto de partida do trabalho}
\scriptsize
\begin{tabularx}{\textwidth}{c >{\bfseries}p{2.7cm}c c c c X}
\toprule
Circ. & Uso & Potência & $I_b$ em 220 V & DJ & Seção inicial & Comando/proteção \\
\midrule
C1 & iluminação & 0,40 kVA & 1,82 A & 10 A & 1,5 mm$^2$ & simples/paralelo/sensor \\
C2 & TUG sala/quarto & 1,20 kVA & 5,45 A & 16 A & 2,5 mm$^2$ & DR \\
C3 & TUG copa & 2,10 kVA & 9,55 A & 20 A & 2,5 mm$^2$ & DR \\
C4 & aquecedor & 5,50 kW & 25,0 A & 32 A 2P & 6 mm$^2$ & DR + dedicado \\
C5 & ar-condicionado & 1,20 kW & 5,93 A & 16 A & 2,5 mm$^2$ & dedicado \\
\bottomrule
\end{tabularx}
\begin{caixaAt}
Valores são hipóteses. Validar placa, método de instalação, temperatura, agrupamento, queda, curto, terminais, DR e especificações dos fabricantes antes de liberar a compra.
\end{caixaAt}
\end{frame}

\begin{frame}{Entrada didática — relacionar carga, padrão e QD}
\small
\begin{columns}[T,onlytextwidth]
\column{0.47\textwidth}
\[
P_{inst}=0{,}4+1{,}2+2{,}1+5{,}5+1{,}2=10{,}4\ \mathrm{kW}.
\]
Na tabela 220/380 V da NT.00001.EQTL Rev.09, uma unidade monofásica maranhense entre 8,1 e 12 kW usa, na referência didática:
\begin{itemize}
\item disjuntor do padrão: 63 A;
\item cobre do cliente: mínimo 10 mm$^2$;
\item demais colunas/notas obrigatórias.
\end{itemize}
\column{0.51\textwidth}
\begin{tikzpicture}[
  b/.style={draw=corPrim,fill=corDef!7,rounded corners,minimum width=2cm,minimum height=.68cm,align=center,font=\scriptsize}]
\node[b](r)at(0,3.6){rede / padrão};
\node[b](m)at(0,2.4){medição + DG};
\node[b](a)at(0,1.2){alimentador};
\node[b](q)at(0,0){QD1 + DPS + DR};
\node[b,draw=corNorma,fill=corNorma!7](pe)at(3,1.2){PE / BEP};
\draw[-{Stealth},thick,corPrim](r)--(m)--(a)--(q);
\draw[corNorma,thick](pe)--(q);\draw[corNorma,thick](pe.south)--(3,-.2)node[ground]{};
\end{tikzpicture}
\begin{caixaAt}
Dimensionar novamente o alimentador interno; não copiar automaticamente a seção do padrão.
\end{caixaAt}
\end{columns}
\end{frame}

\begin{frame}{Lista de condutores por trecho — exemplo de C1}
\scriptsize
\begin{tabularx}{\textwidth}{>{\bfseries}p{2.3cm}p{2.6cm}X X}
\toprule
Trecho & Formação funcional & Identificação & Observação \\
\midrule
QD1--CX1 & F + N + PE & C1-F, C1-N, C1-PE & alimentação da caixa principal \\
CX1--S1p & comum + V1 + V2 & C1-C, C1-V1, C1-V2 & primeiro paralelo \\
CX1--S2p & V1 + V2 + retorno & C1-V1, C1-V2, C1-R1 & segundo paralelo \\
CX1--L1 & retorno + N + PE & C1-R1, C1-N, C1-PE & luminária da sala \\
CX1--S3--L2 & F + retorno / R + N + PE & C1-F/R2/N/PE & copa \\
CX1--sensor--L4 & F + N + saída / R + N + PE & C1-F/N/R4/PE & conferir modelo do sensor \\
\bottomrule
\end{tabularx}
\begin{caixaProjeto}[title=Atividade]
Marcar esses códigos na planta, na lista de cabos, nas anilhas e na folha de continuidade. Qualquer alteração deve aparecer nos quatro lugares.
\end{caixaProjeto}
\end{frame}

\begin{frame}{Unifilar simplificado do QD1}
\centering
\begin{tikzpicture}[
  b/.style={draw=corPrim,fill=corDef!7,rounded corners,minimum width=1.75cm,minimum height=.62cm,align=center,font=\tiny}]
\node[b](dg)at(0,4.0){DG 2P 63 A};
\node[b](dps)at(0,3.0){DPS + barras N/PE};
\node[b,draw=corNorma,fill=corNorma!7](dr)at(0,2.0){DR/RCBO conforme projeto};
\node[b](c1)at(-4.8,.3){C1 luz\\10 A / 1,5};
\node[b](c2)at(-2.4,-.7){C2 TUG\\16 A / 2,5};
\node[b](c3)at(0,.3){C3 copa\\20 A / 2,5};
\node[b](c4)at(2.4,-.7){C4 AQ\\32 A / 6};
\node[b](c5)at(4.8,.3){C5 AC\\16 A / 2,5};
\draw[-{Stealth},thick,corPrim](dg)--(dps)--(dr);
\foreach \x in {c1,c2,c3,c4,c5}{\draw[-{Stealth},thick,corPrim](dr)--(\x);}
\end{tikzpicture}
\begin{caixaObs}
O aluno deve completar: capacidade de interrupção, tipo/polos do DR, ligação do DPS, método/seções finais, correntes corrigidas, quedas, comprimentos e bornes.
\end{caixaObs}
\end{frame}

\begin{frame}{Lista de materiais — comprar por função e requisito}
\scriptsize
\begin{tabularx}{\textwidth}{c >{\bfseries}p{3.0cm}c X}
\toprule
Item & Descrição & Quant. & Informação para seleção \\
\midrule
01 & quadro de distribuição & 1 & módulos, reserva, barras, IP e corrente \\
02 & disjuntores/DPS/DR & conjunto & dados completos do unifilar \\
03 & cabos 1,5/2,5/6/10 mm$^2$ & por medição & classe, isolação, cor e comprimento + reserva \\
04 & eletrodutos e caixas & por rota & diâmetro, ocupação, curvas, fixação e acesso \\
05 & tomadas, interruptores e sensor & por planta & corrente, tensão, módulos, IP e compatibilidade LED \\
06 & luminárias e driver & 4 conjuntos & fluxo, potência, CCT, IRC e montagem \\
07 & terminais/conectores & por borne & seção, classe, corrente e ferramenta \\
08 & etiquetas/anilhas & conjunto & códigos do projeto e resistência ambiental \\
\bottomrule
\end{tabularx}
\begin{caixaAt}
Comprimento de cabo vem da rota executiva, não da distância em linha reta da planta. Incluir subidas, descidas, curvas, quadro, caixas e folgas previstas — sem sobras ocultas excessivas.
\end{caixaAt}
\end{frame}

\begin{frame}{Plano de execução — sequência e pontos de espera}
\scriptsize
\begin{enumerate}
\item revisar documentos, compatibilizar e liberar marcação;
\item marcar caixas, quadro, eixos e cotas;
\item instalar infraestrutura e conferir rotas;
\item \textbf{ponto de espera 1}: inspeção antes de fechar paredes/forro;
\item limpar/provar eletrodutos e lançar cabos identificados;
\item terminar caixas, cargas e QD com torque controlado;
\item \textbf{ponto de espera 2}: inspeção e ensaios desenergizados;
\item corrigir não conformidades e atualizar desenhos;
\item energizar circuito a circuito e testar funções;
\item \textbf{ponto de espera 3}: aceite, as built e entrega.
\end{enumerate}
\begin{caixaObs}
Cada equipe recebe uma frente, mas a passagem de serviço depende do registro da anterior. Isso ensina execução como processo coordenado, não como tarefas soltas.
\end{caixaObs}
\end{frame}

\begin{frame}{Ordem de serviço — exemplo para C1}
\scriptsize
\begin{tabularx}{\textwidth}{>{\bfseries}p{2.4cm}X X}
\toprule
Campo & Conteúdo & Evidência \\
\midrule
Escopo & montar C1: sala paralela, copa simples, banheiro e sensor & planta PE-01 + diagrama DE-01 R00 \\
Fontes & QD1 e qualquer retorno/temporário & lista de bloqueio e ausência de tensão \\
Materiais & cabos, mecanismos, caixas, terminais e luminárias & lote e conferência de recebimento \\
Ferramentas & decapador, crimpador, torquímetro, instrumentos & inspeção e calibração aplicável \\
Conexões & conforme lista de bornes e fios & inspeção ponto a ponto \\
Ensaios & PE, isolação, polaridade e função & folha E-C1 \\
Critério & projeto, procedimento e fabricante & aceite/NC/retrabalho \\
\bottomrule
\end{tabularx}
\begin{caixaAt}
A ordem de serviço não autoriza improvisar alteração. Divergência de campo gera consulta, decisão registrada e revisão correspondente.
\end{caixaAt}
\end{frame}

\begin{frame}{Cartões de falha — aprender diagnóstico sem adivinhação}
\small
\begin{columns}[T,onlytextwidth]
\column{0.49\textwidth}
\begin{caixaFix}[title=Falhas possíveis]
\begin{itemize}
\item comum do paralelo trocado com V1;
\item neutro de C2 na barra do DR errado;
\item PE interrompido em uma tomada;
\item sensor com L e L$'$ invertidos;
\item terminal subcrimpado;
\item etiqueta C4 aplicada no cabo C5;
\item borne do QD sem torque;
\item alteração de rota não documentada.
\end{itemize}
\end{caixaFix}
\column{0.49\textwidth}
\begin{caixaProjeto}[title=Método obrigatório]
\begin{enumerate}
\item definir sintoma;
\item tornar seguro;
\item ler documentos;
\item formular hipótese;
\item escolher ensaio discriminante;
\item corrigir a causa;
\item reensaiar proteção e função;
\item registrar e atualizar as built.
\end{enumerate}
\end{caixaProjeto}
\end{columns}
\end{frame}

\begin{frame}{Ficha de inspeção — aceitar ou rejeitar com evidência}
\scriptsize
\begin{tabularx}{\textwidth}{c X c c X}
\toprule
ID & Verificação & C & NC & Observação/evidência \\
\midrule
I-01 & caixas e cotas correspondem à planta & $\square$ & $\square$ & \\
I-02 & rotas, diâmetros e segregação conferidos & $\square$ & $\square$ & \\
I-03 & cabos, seções e identificações conferem & $\square$ & $\square$ & \\
I-04 & terminais e torques registrados & $\square$ & $\square$ & \\
I-05 & neutros separados por DR e PE contínuo & $\square$ & $\square$ & \\
I-06 & barreiras, tampas e acesso recompostos & $\square$ & $\square$ & \\
I-07 & planta, quadro e unifilar coerentes & $\square$ & $\square$ & \\
I-08 & ensaios anexos atendem aos critérios & $\square$ & $\square$ & \\
\bottomrule
\end{tabularx}
\begin{caixaObs}
“C” exige evidência; campo vazio não significa conforme. “NC” precisa de responsável, prazo, correção e revalidação.
\end{caixaObs}
\end{frame}

\begin{frame}{Dossiê final da unidade didática}
\scriptsize
Cada grupo entrega:
\begin{enumerate}
\item briefing e levantamento da arquitetura;
\item planta elétrica cotada, legenda e detalhes;
\item memória de cargas e dimensionamentos;
\item quadro de cargas e unifilar;
\item esquemas funcional, multifilar e de ligação;
\item lista de cabos, bornes e materiais;
\item análise de risco, ordem de serviço e checklists;
\item registros de torque, inspeção e ensaios;
\item relatório de falha/correção e reensaio;
\item as built e apresentação da instalação funcionando.
\end{enumerate}
\begin{caixaProjeto}[title=Defesa oral]
O avaliador aponta qualquer componente. O aluno deve mostrar onde ele aparece na planta, no quadro, no esquema, na lista de materiais e no relatório de ensaio.
\end{caixaProjeto}
\end{frame}

\begin{frame}{Rubrica — o aluno realmente aprendeu?}
\scriptsize
\begin{tabularx}{\textwidth}{>{\bfseries}p{3.6cm}c X}
\toprule
Competência & Peso & Evidência de domínio \\
\midrule
Leitura/interpretação & 25\% & segue circuitos, referências e bornes sem depender de memorização \\
Projeto/dimensionamento & 25\% & premissas, cálculos e documentos são coerentes e verificáveis \\
Ligação/montagem & 25\% & executa com identificação, ferramenta, terminal, torque e acabamento corretos \\
Segurança/ensaios & 20\% & controla energia, escolhe método, mede, interpreta e registra \\
Comunicação/as built & 5\% & outro profissional entende e mantém a instalação entregue \\
\bottomrule
\end{tabularx}
\begin{caixaAt}
Falha crítica de segurança reprova a prática independentemente da soma. Instalação energizada e funcional não compensa ausência de PE, barreira, proteção ou procedimento.
\end{caixaAt}
\end{frame}

\begin{frame}{Resultado esperado da formação predial}
\centering
\begin{tikzpicture}[
  b/.style={draw=corPrim,fill=corDef!7,rounded corners,minimum width=2.35cm,minimum height=.82cm,align=center,font=\scriptsize},
  a/.style={-{Stealth},very thick,corPrim}]
\node[b](l)at(-4.5,0){Lê\\planta, tabela e diagrama};
\node[b](p)at(-1.5,0){Projeta\\calcula e documenta};
\node[b](e)at(1.5,0){Executa\\liga e identifica};
\node[b](v)at(4.5,0){Verifica\\mede e entrega};
\draw[a](l)--(p)--(e)--(v);
\end{tikzpicture}
\vspace{6mm}
\begin{caixaProjeto}[title=Competência profissional integrada]
O aluno não apenas reproduz um desenho: entende a intenção, confere os dados, transforma o projeto em instalação, prova o resultado e mantém a documentação fiel ao que foi construído.
\end{caixaProjeto}
\end{frame}

%=====================================================================
% REFERENCIAS DO MATERIAL DE LIGACAO, MONTAGEM E VERIFICACAO
%=====================================================================

\section{Referências e continuidade do estudo}

\begin{frame}{Referências normativas e fontes oficiais}
\small
\begin{itemize}
\item \textbf{ABNT NBR 5410} — instalações elétricas de baixa tensão; confirmar edição vigente no \href{https://www.abntcatalogo.com.br/}{Catálogo ABNT}.
\item \textbf{NR-10} — segurança em instalações e serviços em eletricidade: \href{https://www.gov.br/trabalho-e-emprego/pt-br/acesso-a-informacao/participacao-social/conselhos-e-orgaos-colegiados/comissao-tripartite-partitaria-permanente/normas-regulamentadora/normas-regulamentadoras-vigentes/norma-regulamentadora-no-10-nr-10}{página oficial do MTE}.
\item \textbf{Equatorial Energia Maranhão} — \href{https://ma.equatorialenergia.com.br/institucional/leis-e-normas-tecnicas/normas-tecnicas/instalacao-de-padrao/}{normas técnicas e desenhos vigentes de fornecimento}.
\item \textbf{CBMMA} — \href{https://cbm.ssp.ma.gov.br/unidades-bm/unidades-administrativas/dat/legislacao-tecnica-dat/}{legislação técnica de segurança contra incêndio}, incluindo NT 18 para iluminação de emergência.
\item Fichas técnicas, diagramas, torques, compatibilidade e procedimentos dos \textbf{fabricantes efetivamente especificados}.
\item Procedimentos, prontuário, análise de risco e regras de liberação da \textbf{organização responsável pela instalação}.
\end{itemize}
\begin{caixaAt}
Normas e manuais mudam. Registrar número, edição/revisão e data de consulta no memorial e no plano de ensaios.
\end{caixaAt}
\end{frame}

\begin{frame}{Nota temporal sobre a NR-10}
\small
Em agosto de 2026, a página oficial do Ministério do Trabalho e Emprego informa:
\begin{itemize}
\item texto vigente da NR-10 até \textbf{31 de maio de 2027};
\item nova redação aprovada pela \textbf{Portaria MTE nº 737/2026}, com vigência indicada a partir de \textbf{1º de junho de 2027}.
\end{itemize}
\begin{caixaNorma}[title=Aplicação no curso]
Usar o texto legal vigente na data da atividade e preparar a atualização do material/procedimentos antes da mudança de vigência. A existência de uma futura redação não autoriza antecipar ou ignorar obrigações atuais.
\end{caixaNorma}
\begin{caixaAt}
Este Beamer ensina leitura e método, mas não habilita profissional nem substitui capacitação, autorização, análise de risco, supervisão e procedimentos exigidos para o trabalho real.
\end{caixaAt}
\end{frame}




\end{document}
